\documentclass{article}
\usepackage[margin=0.8in]{geometry}
\usepackage{amsmath,amssymb,amsthm,mathtools}
\usepackage{mathrsfs,bm,bbm}
\usepackage{graphicx,booktabs,array,multirow,tabularx,longtable}
\usepackage{enumitem}
\usepackage{xcolor}
\usepackage{algorithm,algpseudocode}
\usepackage{subcaption,tikz}
\usepackage{siunitx,cancel,accents}
\usepackage{microtype}
\usepackage[numbers,sort&compress]{natbib}
\usepackage[colorlinks=true,linkcolor=blue,citecolor=blue,urlcolor=blue]{hyperref}
\usepackage{cleveref}
\setlength{\emergencystretch}{2em}
\newtheorem{assumption}{Assumption}
\newtheorem{definition}{Definition}
\newtheorem{theorem}{Theorem}
\newtheorem{lemma}{Lemma}
\newtheorem{proposition}{Proposition}
\newtheorem{openendedquestion}{Open-ended Question}
\newtheorem{conjecture}{Conjecture}
\newcommand{\coreref}[1]{\ref{#1}}

\begin{document}

\sloppy

\section*{eid\_\allowbreak{}sourcemin\_\allowbreak{}cyclic\_\allowbreak{}effect\_\allowbreak{}rankfrontier}

\noindent\textit{Specialization:} v1\par

\paragraph{TL;DR.} A normalized irreducible overcomplete-ICA law supports an exactly characterizable family of source-minimal cyclic latent completions: its equilibrium-response frontier is the set of deletion-rank-admissible source-column ratios. A shared kernel and matching construction returns every exact same-law witness in worst-case output-optimal \(O(n^3)\) exact arithmetic, conditional on an exact mixing representative. With exact population cumulants, known source count, and all finite source moments, a heterogeneous-order oracle recovers every projective direction after a representation-dependent finite order; moment determinacy is not needed for that positive result. No cutoff or shrinking uniform inference is available on the qualitative class, even on a stricter moment-determinate subclass. A separate common-order iid class with fixed-order moment and separation margins retains simultaneous root-\(N\) confidence unions.

\section{Project justification}

\paragraph{Gap.} Published irreducible OICA results identify projective mixing directions but do not map the postintervention-solvable cyclic completion fiber. Latent-confounded DAG work recursively recovers sources and compatible path matrices under generic higher-cumulant and local latent-count conditions, while fully observed disjoint-cycle work targets graph and coefficient recovery from low-order moment relations. Neither line characterizes the deletion-rank response image of source-minimal latent cyclic completions or constructs every exact same-law intervention witness.

\paragraph{Niche.} The note separates three layers that prior work does not combine: qualitative full-law direction uniqueness; an exact-representative deletion-rank map with constructive cyclic witnesses; and a pointwise heterogeneous-order evaluator under all finite source moments, alongside a distinct quantitatively separated common-order iid inference regime.

\paragraph{Fill.} The deletion-rank frontier gives necessary and sufficient attainability, exact countermodels, representative invariance, and worst-case output-optimal \(O(n^3)\) enumeration. The multi-order handle terminates pointwise with exact cumulants under all finite moments, while explicit near-Gaussian witnesses prove the absence of a uniform cutoff or uniform shrinking inference even on a moment-determinate subclass; the common-order theorem supplies the separate root-\(N\) inference result.

\section{Related work}

DaiAlbrechtSpirtesZhang2026, Section 4, Theorem 2 gives the if-and-only-if graphical criterion for distributional equivalence of irreducible models, Theorem 3 gives traversal by admissible cycle reversals and edge additions or deletions, and Theorem 4 gives the maximal-class presentation. Those graphical equivalence, traversal, and maximal-presentation results do not enumerate the fixed-law postintervention response image, prove deletion-rank equality, or materialize exact same-law effect witnesses. The present response theorem deliberately uses a source-minimal, all-non-Gaussian exact-law completion fiber, a scope narrower than Dai et al.'s graphically irreducible model family. TramontanoKivvaSalehkaleybarDrtonKiyavash2024, Theorem 3.3 gives an if-and-only-if criterion for generic identification of the total causal effect from \(j\) to \(i\) without graph knowledge. That acyclic selected-effect criterion does not enumerate every response over a cyclic exact same-law completion fiber, prove deletion-rank sharpness, or construct same-law countermodels. TramontanoKivvaSalehkaleybarDrtonKiyavash2025, Theorem 3.5 states that the \(T\)-to-\(Y\) causal effect is generically identifiable from the first \(k(\ell)\) cumulants of the observational distribution. That restricted acyclic higher-cumulant selected-effect result does not characterize all effects across exact same-law cyclic completions, the fixed-law postintervention response image, deletion-rank equality, or exact countermodel witnesses. TramontanoDrtonEtesami2026 studies coefficient identification for a known mixed graph and obtains only partial cyclic criteria, a different conditioning object from the fixed source-direction fiber. In LacerdaSpirtesRamseyHoyer2008, Section 4 constructs square candidates by admissible row permutations, while Theorem 1 states that LiNG-D returns a distribution-entailment equivalence class. The equality of the candidates' equilibrium-effect projection with \(\mathcal R_{xy}(C)\) is the new square-case reduction here. BongersForrePetersMooij2021 supplies the distinction between observational and interventional solvability. WangSeigal2026 supplies the OICA identification and coupled-decomposition benchmark, AnandkumarGeJanzamin2015 supplies the tensor perturbation technique, and ChenBickel2006 supplies the asymptotic ICA motivation; none studies a rank-selected causal set. PodosinnikovaEtAl2019's Algorithm 1 estimates a cumulant-generating-function-Hessian subspace and applies an SDP relaxation; its model-specific high-probability component-recovery guarantee for uniformly spherical mixing components with \(k<(2-\epsilon)p\log p\) is not challenged here. The heterogeneous-order boundary instead rules out a different proposed theorem: a predeclared finite-cumulant-cutoff, uniformly consistent direction/frontier or uniformly shrinking confidence-union extension over the unseparated qualitative all-finite-moments class. VirtaNordhausenOja2016 remains a square third/fourth-cumulant comparison. NandyMaathuisRichardson2017 confirms that joint interventions require globally compatible assignments rather than Cartesian products, motivating but not subsuming the deferred joint-frontier seed.

\section{Setup}

Target (estimand / structural parameter): \(\boldsymbol\Theta_x^{\mathrm{law}}(P_X)=\{\rho_x(\widetilde{\mathcal M}):\widetilde{\mathcal M}\in\mathfrak F_x^{\mathrm{law}}(P_X)\}\), with ordered scalar target \(\Theta_{y\leftarrow x}^{\mathrm{law}}(P_X)=\{[\rho_x(\widetilde{\mathcal M})]_y:\widetilde{\mathcal M}\in\mathfrak F_x^{\mathrm{law}}(P_X)\}\).
Primary functional / estimator / diagnostic: Full-law uniqueness identifies \(\mathcal D(P_X)=\{[c_j]:j\in\mathsf S\}\) set-theoretically. For any exact irreducible mixing representative \(C\) of this projective set,
\[
 \boldsymbol\Theta_x^{\mathrm{law}}(P_X)
 =\mathcal V_x(C)
 =\left\{\left(\frac{C_{ij}}{C_{xj}}\right)_{i\in\mathsf O\setminus\{x\}}:j\in J_x(C)\right\},
 \qquad
 \Theta_{y\leftarrow x}^{\mathrm{law}}(P_X)
 =\mathcal R_{xy}(C)
 =\left\{\frac{C_{yj}}{C_{xj}}:j\in J_x(C)\right\}.
\]
These expressions are independent of the representative. Conditional on exact \(C\), the shared-kernel and shared-matching certificate evaluates all frontier values and materializes one exact dense witness per value in worst-case output-optimal \(O(n^3)\) exact arithmetic; this is not a procedure or complexity bound for recovering \(C\) from \(P_X\). The qualitative full-law result does not evaluate \(\mathcal D(P_X)\). Under the heterogeneous theorem's full assumptions---known source count \(n\), normalized full-row-rank pairwise-nonproportional mixing, mutually independent nondegenerate non-Gaussian sources, and finite absolute moments of every positive integer order---the multi-order handle evaluates the direction set by a terminating pointwise population-oracle procedure. Moment determinacy is not used for this positive recovery. The negative nonuniformity conclusion holds even on the stricter moment-determinate subclass. The population-oracle result is incomparable with the accepted lower-order bivariate arrow results because it recovers all projective directions for an arbitrary observed dimension and then maps them to a cyclic response frontier, whereas those results identify a single axis-conditioned bivariate arrow. Under the separated common-order iid assumptions, which retain their fixed-order moment envelope, the declared confidence unions provide simultaneous root-\(N\) inference..
\begin{itemize}
  \item \texttt{p} --- \texttt{observed dimension}; \(\{2,3,\ldots\}\); \texttt{fixed integer}
  \item \texttt{n} --- \texttt{source and completion dimension}; \(\{p,p+1,\ldots\}\); \texttt{fixed integer}
  \item \texttt{N} --- \texttt{sample size}; \(\mathbb N\); \texttt{inference index}
  \item \texttt{d} --- \texttt{tensor lifting degree}; \(\{1,2,\ldots\}\); \texttt{fixed integer}
  \item \texttt{q} --- \texttt{tensor contraction degree}; \(\{1,2,\ldots\}\); \texttt{fixed integer}
  \item \texttt{r} --- \texttt{common cumulant order}; \(\{3,4,\ldots\}\)
  \par\smallskip\noindent\emph{Definition.} \(r=2d+q\)
  \item \texttt{alpha} --- \texttt{miscoverage level}; \((0,1)\); \texttt{fixed constant}
  \item \texttt{a} --- \texttt{treatment-\allowbreak{}loading margin}; \((0,1]\); \texttt{fixed constant}
  \item \texttt{delta} --- \texttt{matrix and deletion-\allowbreak{}minor margin}; \((0,1]\); \texttt{fixed constant}
  \item \texttt{sigma} --- \texttt{tensor and probe separation margin}; \((0,1]\); \texttt{fixed constant}
  \item \texttt{kappa} --- \texttt{lower cumulant magnitude}; \((0,\infty)\); \texttt{fixed constant}
  \item \texttt{Lambda} --- \texttt{upper cumulant magnitude}; \([\kappa,\infty)\); \texttt{fixed constant}
  \item \texttt{M} --- \texttt{source moment envelope}; \([1,\infty)\); \texttt{fixed constant}
  \item \texttt{Obs} --- \texttt{observed index set}; \(2^{\{1,\ldots,n\}}\)
  \par\smallskip\noindent\emph{Definition.} \(\mathsf O=\{1,\ldots,p\}\)
  \item \texttt{Src} --- \texttt{source index set}; \(2^{\{1,\ldots,n\}}\)
  \par\smallskip\noindent\emph{Definition.} \(\mathsf S=\{1,\ldots,n\}\)
  \item \texttt{Lat} --- \texttt{latent completion index set}; \(2^{\{1,\ldots,n\}}\)
  \par\smallskip\noindent\emph{Definition.} \(\mathsf L=\{p+1,\ldots,n\}\)
  \item \texttt{x} --- \texttt{treatment coordinate}; \(\mathsf O\)
  \item \texttt{y} --- \texttt{outcome coordinate}; \(\mathsf O\setminus\{x\}\)
  \item \texttt{epsilon} --- \texttt{source vector}; \(\mathbb R^n\)
  \par\smallskip\noindent\emph{Definition.} \(\varepsilon=(\varepsilon_1,\ldots,\varepsilon_n)^\top\)
  \item \texttt{C} --- \texttt{normalized overcomplete mixing matrix}; \(\mathbb R^{p\times n}\)
  \par\smallskip\noindent\emph{Definition.} \(C=[c_1\;\cdots\;c_n]\)
  \item \texttt{X} --- \texttt{observed random vector}; \(\mathbb R^p\)
  \par\smallskip\noindent\emph{Definition.} \(X=C\varepsilon\)
  \item \texttt{PX} --- \texttt{observational law}; \(\mathcal P(\mathbb R^p)\)
  \par\smallskip\noindent\emph{Definition.} \(P_X=\mathcal L(C\varepsilon)\)
  \item \texttt{V} --- \texttt{endogenous completion vector}; \(\mathbb R^n\)
  \item \texttt{B} --- \texttt{structural coefficient matrix}; \(\mathbb R^{n\times n}\)
  \item \texttt{Qmat} --- \texttt{structural equation matrix}; \(\mathbb R^{n\times n}\)
  \par\smallskip\noindent\emph{Definition.} \(Q=I_n-B\)
  \item \texttt{Hmat} --- \texttt{source-\allowbreak{}assignment matrix}; \(\mathbb R^{n\times n}\)
  \item \texttt{Mcomp} --- \texttt{source-\allowbreak{}minimal cyclic structural completion}
  \par\smallskip\noindent\emph{Definition.} \(\mathcal M=(Q,H)\)
  \item \texttt{Fx} --- \texttt{postintervention-\allowbreak{}solvable fixed-\allowbreak{}law completion fiber}
  \item \texttt{sassign} --- \texttt{equation-\allowbreak{}to-\allowbreak{}source assignment}; \(\mathsf O\to\mathsf S\)
  \par\smallskip\noindent\emph{Definition.} \(s_{\mathcal M}(i)\) is the unique \(j\) for which \(H_{ij}\ne0\)
  \item \texttt{theta} --- \texttt{hard-\allowbreak{}intervention equilibrium effect}; \(\mathbb R\)
  \par\smallskip\noindent\emph{Definition.} \(\theta_{y\leftarrow x}(\mathcal M)=(Q^{-1})_{yx}/(Q^{-1})_{xx}\)
  \item \texttt{Jx} --- \texttt{deletion-\allowbreak{}rank-\allowbreak{}admissible source set}; \(2^{\mathsf S}\)
  \par\smallskip\noindent\emph{Definition.} \(J_x(C)=\{j\in\mathsf S:C_{xj}\ne0,\ \operatorname{rank}(C_{-x,-j})=p-1\}\)
  \item \texttt{Rxy} --- \texttt{deletion-\allowbreak{}rank effect frontier}; \(2^{\mathbb R}\)
  \par\smallskip\noindent\emph{Definition.} \(\mathcal R_{xy}(C)=\{C_{yj}/C_{xj}:j\in J_x(C)\}\)
  \item \texttt{T0} --- \texttt{basis extension of the observed mixing rows}; \(\operatorname{GL}_n(\mathbb R)\)
  \par\smallskip\noindent\emph{Definition.} \((T_0)_{\mathsf O,:}=C\)
  \item \texttt{W0} --- \texttt{inverse basis extension}; \(\operatorname{GL}_n(\mathbb R)\)
  \par\smallskip\noindent\emph{Definition.} \(W_0=T_0^{-1}\)
  \item \texttt{Cert} --- \texttt{all-\allowbreak{}witness completion algorithm}; \((C,x)\mapsto\{\mathcal M_j:j\in J_x(C)\}\)
  \item \texttt{u} --- \texttt{normalizing tensor probe}; \(\mathbb S^{p-1}\); \texttt{fixed vector}
  \item \texttt{v} --- \texttt{separating tensor probe}; \(\mathbb S^{p-1}\); \texttt{fixed vector}
  \item \texttt{lambda} --- \texttt{source cumulant vector}; \(\mathbb R^n\)
  \par\smallskip\noindent\emph{Definition.} \(\lambda_j=\operatorname{cum}_r(\varepsilon_j)\)
  \item \texttt{zeta} --- \texttt{representation-\allowbreak{}indexed inference parameter}
  \par\smallskip\noindent\emph{Definition.} \(\zeta=(P_X,C,\varepsilon,\lambda)\)
  \item \texttt{Tr} --- order-\(r\) observational cumulant tensor; \(\operatorname{Sym}^r(\mathbb R^p)\)
  \par\smallskip\noindent\emph{Definition.} \(\mathcal T_r=\operatorname{cum}_r(X)=\sum_{j=1}^n\lambda_j c_j^{\otimes r}\)
  \item \texttt{Vd} --- \texttt{lifted direction matrix}; \(\mathbb R^{p^d\times n}\)
  \par\smallskip\noindent\emph{Definition.} \(V_d=[c_1^{\otimes d}\;\cdots\;c_n^{\otimes d}]\)
  \item \texttt{Mreg} --- \texttt{separated common-\allowbreak{}order OICA law class}
  \item \texttt{OICAInv} --- \texttt{tensor-\allowbreak{}pencil inverse}; \(\mathcal T_r\mapsto\{c_1,\ldots,c_n\}\)
  \item \texttt{eta} --- \texttt{contracted tensor singular-\allowbreak{}value bound}; \((0,\infty)\)
  \par\smallskip\noindent\emph{Definition.} \(\eta=\kappa\sigma^{q+2}\)
  \item \texttt{chi} --- \texttt{eigenvector condition bound}; \((0,\infty)\)
  \par\smallskip\noindent\emph{Definition.} \(\chi=\sqrt n/\sigma\)
  \item \texttt{hconst} --- \texttt{pencil perturbation coefficient}; \((0,\infty)\)
  \par\smallskip\noindent\emph{Definition.} \(h=2/\eta+2n\Lambda/\eta^2\)
  \item \texttt{eps0} --- \texttt{local inversion radius}; \((0,\infty)\)
  \par\smallskip\noindent\emph{Definition.} \(\epsilon_0=\min\{\eta/2,\sigma/(6\chi h)\}\)
  \item \texttt{Lz} --- \texttt{ratio-\allowbreak{}coordinate inverse modulus}; \((0,\infty)\)
  \par\smallskip\noindent\emph{Definition.} \(L_z=nh\{2\chi+6n\Lambda\chi^2/(\eta\sigma)\}\)
  \item \texttt{LC} --- \texttt{mixing inverse modulus}; \((0,\infty)\)
  \par\smallskip\noindent\emph{Definition.} \(L_C=2\sqrt{np}L_z\)
  \item \texttt{Xsample} --- \texttt{iid observational sample}; \((\mathbb R^p)^N\)
  \par\smallskip\noindent\emph{Definition.} \(\mathbf X_N=(X^{(1)},\ldots,X^{(N)})\)
  \item \texttt{momentdim} --- number of nonconstant moments through order \(r\); \(\mathbb N\)
  \par\smallskip\noindent\emph{Definition.} \(m=\binom{p+r}{r}-1\)
  \item \texttt{Hbound} --- \texttt{raw-\allowbreak{}moment second-\allowbreak{}moment envelope}; \([1,\infty)\)
  \par\smallskip\noindent\emph{Definition.} \(H=n^{2r}\max\{1,M\}\)
  \item \texttt{Bmom} --- \texttt{moment clipping level}; \([2,\infty)\)
  \par\smallskip\noindent\emph{Definition.} \(B_m=1+H\)
  \item \texttt{eN} --- \texttt{simultaneous raw-\allowbreak{}moment error radius}; \((0,\infty)\)
  \par\smallskip\noindent\emph{Definition.} \(e_N=\sqrt{m(1+H)/(N\alpha)}\)
  \item \texttt{Lmom} --- \texttt{moment-\allowbreak{}to-\allowbreak{}cumulant Lipschitz coefficient}; \((0,\infty)\)
  \par\smallskip\noindent\emph{Definition.} \(L_m=p^{r/2}\sum_{\pi\in\mathfrak P_r}(|\pi|-1)!|\pi|B_m^{|\pi|-1}\), where \(\mathfrak P_r\) is the set of partitions of \(\{1,\ldots,r\}\)
  \item \texttt{tauN} --- \texttt{empirical cumulant-\allowbreak{}tensor error radius}; \((0,\infty)\)
  \par\smallskip\noindent\emph{Definition.} \(\tau_N=L_m e_N\)
  \item \texttt{That} --- \texttt{clipped empirical cumulant tensor}; \(\operatorname{Sym}^r(\mathbb R^p)\)
  \item \texttt{Chat} --- \texttt{minimum-\allowbreak{}distance normalized mixing estimate}; \(\mathbb R^{p\times n}\)
  \item \texttt{Jhat} --- \texttt{estimated deletion-\allowbreak{}rank-\allowbreak{}admissible source set}; \(2^{\mathsf S}\)
  \par\smallskip\noindent\emph{Definition.} \(\widehat J_x=J_x(\widehat C)\)
  \item \texttt{bN} --- \texttt{mixing-\allowbreak{}matrix confidence radius}; \((0,\infty)\)
  \par\smallskip\noindent\emph{Definition.} \(b_N=2L_C\tau_N\)
  \item \texttt{wN} --- \texttt{effect-\allowbreak{}coordinate confidence radius}; \((0,\infty)\)
  \par\smallskip\noindent\emph{Definition.} \(w_N=(a^{-1}+a^{-2})b_N\)
  \item \texttt{Uxy} --- \texttt{simultaneous confidence union for the effect frontier}; \(2^{\mathbb R}\)
  \item \texttt{Hmulti} --- \texttt{heterogeneous-\allowbreak{}order recovery handle}
  \item \texttt{Mlaw} --- \texttt{observational-\allowbreak{}law structural completion}; \texttt{variable}
  \item \texttt{mMlaw} --- \texttt{source and equation count of an observational-\allowbreak{}law completion}; \(\{p,p+1,\ldots\}\)
  \par\smallskip\noindent\emph{Definition.} \(m_{\widetilde{\mathcal M}}\)
  \item \texttt{Qlaw} --- \texttt{law-\allowbreak{}completion structural equation matrix}; \(\mathbb R^{m_{\widetilde{\mathcal M}}\times m_{\widetilde{\mathcal M}}}\)
  \par\smallskip\noindent\emph{Definition.} \(Q^{\widetilde{\mathcal M}}\)
  \item \texttt{Hlaw} --- \texttt{law-\allowbreak{}completion source-\allowbreak{}assignment matrix}; \(\mathbb R^{m_{\widetilde{\mathcal M}}\times m_{\widetilde{\mathcal M}}}\)
  \par\smallskip\noindent\emph{Definition.} \(H^{\widetilde{\mathcal M}}\)
  \item \texttt{xilaw} --- \texttt{law-\allowbreak{}completion source vector}; \(\mathbb R^{m_{\widetilde{\mathcal M}}}\)
  \par\smallskip\noindent\emph{Definition.} \(\xi^{\widetilde{\mathcal M}}\)
  \item \texttt{Vlaw} --- \texttt{law-\allowbreak{}completion endogenous equilibrium vector}; \(\mathbb R^{m_{\widetilde{\mathcal M}}}\)
  \par\smallskip\noindent\emph{Definition.} \(V^{\widetilde{\mathcal M}}\)
  \item \texttt{Alaw} --- \texttt{observed mixing matrix of a law completion}; \(\mathbb R^{p\times m_{\widetilde{\mathcal M}}}\)
  \par\smallskip\noindent\emph{Definition.} \(A^{\widetilde{\mathcal M}}=[(Q^{\widetilde{\mathcal M}})^{-1}H^{\widetilde{\mathcal M}}]_{\mathsf O,:}\)
  \item \texttt{Flaw} --- \texttt{source-\allowbreak{}minimal observational-\allowbreak{}law completion fiber}
  \item \texttt{slaw} --- \texttt{law-\allowbreak{}completion equation-\allowbreak{}to-\allowbreak{}source assignment}; \(\mathsf O\to\{1,\ldots,m_{\widetilde{\mathcal M}}\}\)
  \par\smallskip\noindent\emph{Definition.} \(s_{\widetilde{\mathcal M}}^{\mathrm{law}}(i)\) is the unique \(k\) with \(H^{\widetilde{\mathcal M}}_{ik}\ne0\)
  \item \texttt{salign} --- \texttt{reference-\allowbreak{}direction-\allowbreak{}aligned equation source}; \(\mathsf O\to\mathsf S\)
  \par\smallskip\noindent\emph{Definition.} \(\bar s_{\widetilde{\mathcal M}}(i)\)
  \item \texttt{rhoLaw} --- \texttt{observed hard-\allowbreak{}intervention equilibrium response vector}; \(\mathbb R^{p-1}\)
  \par\smallskip\noindent\emph{Definition.} \(\rho_x(\widetilde{\mathcal M})\)
  \item \texttt{Vfront} --- \texttt{deletion-\allowbreak{}rank response-\allowbreak{}vector frontier}; \(2^{\mathbb R^{p-1}}\)
  \par\smallskip\noindent\emph{Definition.} \(\mathcal V_x(C)\)
  \item \texttt{wvecN} --- \texttt{response-\allowbreak{}vector confidence radius}; \((0,\infty)\)
  \par\smallskip\noindent\emph{Definition.} \(w_N^{\mathrm{vec}}=\sqrt{p-1}\,w_N\)
  \item \texttt{Uvec} --- \texttt{simultaneous confidence union for the response-\allowbreak{}vector frontier}; \(2^{\mathbb R^{p-1}}\)
  \item \texttt{Dlaw} --- \texttt{observational-\allowbreak{}law projective direction functional}; \(2^{\mathbb P^{p-1}}\)
  \par\smallskip\noindent\emph{Definition.} \(\mathcal D(P_X)\)
\end{itemize}

\section{Assumptions}

\begin{assumption}[ass:full-row-rank]\label{ass:full-row-rank}
\(\operatorname{rank}(C)=p\)
\par\smallskip\noindent\emph{Standard} (full-row-rank overcomplete mixing, \cite{WangSeigal2026}).
\end{assumption}

\begin{assumption}[ass:nonzero-columns]\label{ass:nonzero-columns}
\(c_j\ne0\) for every \(j\in\mathsf S\)
\par\smallskip\noindent\emph{Standard} (nondegenerate ICA directions, \cite{WangSeigal2026}).
\end{assumption}

\begin{assumption}[ass:distinct-directions]\label{ass:distinct-directions}
\(c_j\notin\operatorname{span}(c_k)\) for all distinct \(j,k\in\mathsf S\)
\par\smallskip\noindent\emph{Standard} (no proportional ICA components, \cite{WangSeigal2026}).
\end{assumption}

\begin{assumption}[ass:unit-normalization]\label{ass:unit-normalization}
\(\lVert c_j\rVert_2=1\) for every \(j\in\mathsf S\)
\par\smallskip\noindent\emph{Standard} (ICA scale normalization, \cite{Comon1994}).
\end{assumption}

\begin{assumption}[ass:independent-sources]\label{ass:independent-sources}
\(\varepsilon_1,\ldots,\varepsilon_n\) are mutually independent
\par\smallskip\noindent\emph{Standard} (independent component sources, \cite{Comon1994}).
\end{assumption}

\begin{assumption}[ass:nondegenerate-sources]\label{ass:nondegenerate-sources}
\(\operatorname{Var}(\varepsilon_j)>0\) for every \(j\in\mathsf S\)
\par\smallskip\noindent\emph{Standard} (nondegenerate independent components, \cite{Comon1994}).
\end{assumption}

\begin{assumption}[ass:nongaussian-sources]\label{ass:nongaussian-sources}
\(\varepsilon_j\) is non-Gaussian for every \(j\in\mathsf S\)
\par\smallskip\noindent\emph{Standard} (non-Gaussian independent components, \cite{Comon1994}).
\end{assumption}

\begin{assumption}[ass:moment-determinate-sources]\label{ass:moment-determinate-sources}
For every source \(j\in\mathsf S\) and every integer \(k\ge1\), \(\mathbb E|\varepsilon_j|^k<\infty\)
\par\smallskip\noindent\emph{Standard} (all finite source moments; the stable internal identifier is retained only for atomic compatibility, and the mathematical premise is all finite moments, not moment determinacy, \cite{Marcinkiewicz1939}).
\end{assumption}

\begin{assumption}[ass:unit-structural-diagonal]\label{ass:unit-structural-diagonal}
\(Q_{ii}=1\) for every \(i\in\mathsf S\)
\par\smallskip\noindent\emph{Standard} (zero structural self-loop normalization, \cite{LacerdaSpirtesRamseyHoyer2008}).
\end{assumption}

\begin{assumption}[ass:source-assignment]\label{ass:source-assignment}
\(H\) is monomial, meaning that every row and every column has exactly one nonzero entry
\par\smallskip\noindent\emph{Standard} (one independent disturbance per LiNG vertex, up to source permutation and nonzero scaling, \cite{DaiAlbrechtSpirtesZhang2026}).
\end{assumption}

\begin{assumption}[ass:observational-solvability]\label{ass:observational-solvability}
\(\det(Q)\ne0\)
\par\smallskip\noindent\emph{Standard} (unique solvability of a linear cyclic SCM, \cite{BongersForrePetersMooij2021}).
\end{assumption}

\begin{assumption}[ass:fixed-law-match]\label{ass:fixed-law-match}
\([Q^{-1}H]_{\mathsf O,:}=C\)
\par\smallskip\noindent\emph{Novel.} This is the fixed-fiber conditioning convention: one supplied exact irreducible mixing representative is held fixed for the structural witness certificate while latent structural rows vary. It does not assert that qualitative full-law uniqueness evaluates or recovers that representative.
\end{assumption}

\begin{assumption}[ass:postintervention-solvability]\label{ass:postintervention-solvability}
\(\det(Q_{-x,-x})\ne0\)
\par\smallskip\noindent\emph{Standard} (unique solvability after a perfect intervention, \cite{BongersForrePetersMooij2021}).
\end{assumption}

\begin{assumption}[ass:iid-sampling]\label{ass:iid-sampling}
\(X^{(1)},\ldots,X^{(N)}\) are independent draws from \(P_X\)
\par\smallskip\noindent\emph{Standard} (iid ICA sampling, \cite{ChenBickel2006}).
\end{assumption}

\begin{assumption}[ass:cumulant-magnitude-margin]\label{ass:cumulant-magnitude-margin}
\(\kappa\le |\lambda_j|\le\Lambda\) for every \(j\in\mathsf S\)
\par\smallskip\noindent\emph{Standard} (nonvanishing common-order cumulant margin, \cite{AnandkumarGeJanzamin2015}).
\end{assumption}

\begin{assumption}[ass:probe-orientation-margin]\label{ass:probe-orientation-margin}
\(u^\top c_j\ge\sigma\) for every \(j\in\mathsf S\)
\par\smallskip\noindent\emph{Standard} (oriented tensor-probe margin, \cite{AnandkumarGeJanzamin2015}).
\end{assumption}

\begin{assumption}[ass:lifted-rank-margin]\label{ass:lifted-rank-margin}
\(s_{\min}(V_d)\ge\sigma\)
\par\smallskip\noindent\emph{Standard} (robust lifted-factor rank, \cite{WangSeigal2026}).
\end{assumption}

\begin{assumption}[ass:pencil-gap]\label{ass:pencil-gap}
\(\left|v^\top c_j/(u^\top c_j)-v^\top c_k/(u^\top c_k)\right|\ge\sigma\) for all distinct \(j,k\in\mathsf S\)
\par\smallskip\noindent\emph{Standard} (simple-spectrum tensor-pencil separation, \cite{AnandkumarGeJanzamin2015}).
\end{assumption}

\begin{assumption}[ass:source-moment-envelope]\label{ass:source-moment-envelope}
\(\mathbb E|\varepsilon_j|^{2r}\le M\) for every \(j\in\mathsf S\)
\par\smallskip\noindent\emph{Standard} (fixed-order moment envelope, \cite{AnandkumarGeJanzamin2015}).
\end{assumption}

\begin{assumption}[ass:mixing-singular-margin]\label{ass:mixing-singular-margin}
\(s_{\min}(C)\ge\delta\)
\par\smallskip\noindent\emph{Standard} (uniform full-row-rank margin, \cite{WangSeigal2026}).
\end{assumption}

\begin{assumption}[ass:loading-selection-margin]\label{ass:loading-selection-margin}
\(|C_{ij}|=0\) or \( |C_{ij}|\ge a\) for every \(i\in\mathsf O\) and \(j\in\mathsf S\)
\par\smallskip\noindent\emph{Standard} (beta-min loading separation, \cite{AnandkumarGeJanzamin2015}).
\end{assumption}

\begin{assumption}[ass:deletion-rank-margin]\label{ass:deletion-rank-margin}
\(s_{\min}(C_{-i,-j})=0\) or \(s_{\min}(C_{-i,-j})\ge\delta\) for every \(i\in\mathsf O\) and \(j\in\mathsf S\)
\par\smallskip\noindent\emph{Standard} (rank-selection singular-value margin, \cite{WangSeigal2026}).
\end{assumption}

\begin{assumption}[ass:regular-class-nonempty]\label{ass:regular-class-nonempty}
\(\mathfrak P_{p,n,r}(a,\delta,\sigma,\kappa,\Lambda,M)\ne\varnothing\)
\par\smallskip\noindent\emph{Novel.} The separated inference regime is restricted to dimension and margin choices attained by at least one normalized OICA tuple. Such regimes include generic collections of at most \(\binom{p+d-1}{d}\) lifted directions with nonzero common-order cumulants and probes chosen away from their finitely many collision hyperplanes.
\end{assumption}

\section{Definitions}

\begin{definition}[def:completion-fiber]\label{def:completion-fiber}
\par\smallskip\noindent\emph{Defined object.} \texttt{Fixed-\allowbreak{}law source-\allowbreak{}minimal completion fiber}
\par\smallskip\noindent\emph{Construction.} \(\mathfrak F_x(C,\varepsilon)=\{\mathcal M=(Q,H):QV=H\varepsilon,\ Q,H\in\mathbb R^{n\times n},\ Q_{ii}=1\ (i\in\mathsf S),\ H\text{ is monomial},\ \det(Q)\ne0,\ [Q^{-1}H]_{\mathsf O,:}=C,\ \det(Q_{-x,-x})\ne0\}\)
\par\smallskip\noindent\emph{Member properties.} \texttt{ass:\allowbreak{}unit-\allowbreak{}structural-\allowbreak{}diagonal}; \texttt{ass:\allowbreak{}source-\allowbreak{}assignment}; \texttt{ass:\allowbreak{}observational-\allowbreak{}solvability}; \texttt{ass:\allowbreak{}fixed-\allowbreak{}law-\allowbreak{}match}; \texttt{ass:\allowbreak{}postintervention-\allowbreak{}solvability}.
\end{definition}

\begin{definition}[def:assigned-source]\label{def:assigned-source}
\par\smallskip\noindent\emph{Defined object.} \texttt{Assigned source}
\par\smallskip\noindent\emph{Construction.} For \(\mathcal M=(Q,H)\in\mathfrak F_x(C,\varepsilon)\), set \(s_{\mathcal M}(i)=j\) exactly when \(H_{ij}\ne0\)
\par\smallskip\noindent\emph{Inputs.} \(\mathcal M\); \(i\).
\end{definition}

\begin{definition}[def:equilibrium-effect]\label{def:equilibrium-effect}
\par\smallskip\noindent\emph{Defined object.} \texttt{Hard-\allowbreak{}intervention equilibrium effect}
\par\smallskip\noindent\emph{Construction.} \(\theta_{y\leftarrow x}(\mathcal M)=\partial\,\operatorname{Eq}_{\mathcal M}(V_y\mid\operatorname{do}(V_x=t))/\partial t=(Q^{-1})_{yx}/(Q^{-1})_{xx}\)
\par\smallskip\noindent\emph{Inputs.} \(\mathcal M\); \(x\); \(y\).
\end{definition}

\begin{definition}[def:admissible-source-set]\label{def:admissible-source-set}
\par\smallskip\noindent\emph{Defined object.} \texttt{Deletion-\allowbreak{}rank-\allowbreak{}admissible sources}
\par\smallskip\noindent\emph{Construction.} \(J_x(C)=\{j\in\mathsf S:C_{xj}\ne0\text{ and }\operatorname{rank}(C_{-x,-j})=p-1\}\)
\par\smallskip\noindent\emph{Inputs.} \(C\); \(x\).
\end{definition}

\begin{definition}[def:effect-frontier]\label{def:effect-frontier}
\par\smallskip\noindent\emph{Defined object.} \texttt{Deletion-\allowbreak{}rank equilibrium-\allowbreak{}effect frontier}
\par\smallskip\noindent\emph{Construction.} \(\mathcal R_{xy}(C)=\{C_{yj}/C_{xj}:j\in J_x(C)\}\)
\par\smallskip\noindent\emph{Inputs.} \(C\); \(x\); \(y\).
\end{definition}

\begin{definition}[def:witness-algorithm]\label{def:witness-algorithm}
\par\smallskip\noindent\emph{Defined object.} \texttt{Shear-\allowbreak{}and-\allowbreak{}matching completion certificate}
\par\smallskip\noindent\emph{Construction.} Extend the rows of \(C\) to \(T_0\in\operatorname{GL}_n(\mathbb R)\), set \(W_0=T_0^{-1}\), and put \(K_x=(W_0)_{:,\{x\}\cup\mathsf L}\). Determine \(J_x(C)\) by retaining exactly the indices \(j\) for which \(C_{xj}\ne0\) and the \(j\)-th row of \(K_x\) is nonzero. Form the bipartite nonzero-pattern graph of \((W_0)_{:, -x}\), with row vertices \(\mathsf S\) and column vertices \(\mathsf S\setminus\{x\}\). Compute once a matching that saturates all column vertices, its unique unmatched row, and an alternating-reachability forest rooted at that row. For each \(j\in J_x(C)\), flip the matching along a stored alternating path to \(j\), obtaining a matching that saturates all columns other than \(x\) and leaves row \(j\) unmatched. Retain \(T_j=T_0\) and \(W_j=W_0\) if \((W_0)_{jx}\ne0\); otherwise choose \(\ell\in\mathsf L\) with \((W_0)_{j\ell}\ne0\) and set \(T_j=(I_n+e_\ell e_x^\top)T_0\) and \(W_j=W_0(I_n-e_\ell e_x^\top)\). Augment the matching by assigning row \(j\) to column \(x\), and let \(\Pi_j\) be the row-permutation matrix whose row indexed by a column selects the row assigned to that column. With \(D_j=\operatorname{diag}(\Pi_jW_j)\), return \(Q_j=D_j^{-1}\Pi_jW_j\), \(H_j=D_j^{-1}\Pi_j\), and \(\mathcal M_j=(Q_j,H_j)\).
\par\smallskip\noindent\emph{Inputs.} \(C\); \(x\).
\end{definition}

\begin{definition}[def:regular-oica-class]\label{def:regular-oica-class}
\par\smallskip\noindent\emph{Defined object.} \texttt{Separated common-\allowbreak{}order OICA class}
\par\smallskip\noindent\emph{Construction.} \(\mathfrak P_{p,n,r}(a,\delta,\sigma,\kappa,\Lambda,M)=\{\zeta=(P_X,C,\varepsilon,\lambda):X=C\varepsilon\in\mathbb R^p,\ C\in\mathbb R^{p\times n},\lambda_j=\operatorname{cum}_r(\varepsilon_j),\ d\ge1,\ q\ge1,\ r=2d+q,\ n\le\binom{p+d-1}{d},\text{ and the listed source, normalization, cumulant, tensor, moment, loading, and deletion-margin properties hold for the indices }p,n,r,d,q\}\)
\par\smallskip\noindent\emph{Member properties.} \texttt{ass:\allowbreak{}full-\allowbreak{}row-\allowbreak{}rank}; \texttt{ass:\allowbreak{}nonzero-\allowbreak{}columns}; \texttt{ass:\allowbreak{}distinct-\allowbreak{}directions}; \texttt{ass:\allowbreak{}unit-\allowbreak{}normalization}; \texttt{ass:\allowbreak{}independent-\allowbreak{}sources}; \texttt{ass:\allowbreak{}nondegenerate-\allowbreak{}sources}; \texttt{ass:\allowbreak{}nongaussian-\allowbreak{}sources}; \texttt{ass:\allowbreak{}cumulant-\allowbreak{}magnitude-\allowbreak{}margin}; \texttt{ass:\allowbreak{}probe-\allowbreak{}orientation-\allowbreak{}margin}; \texttt{ass:\allowbreak{}lifted-\allowbreak{}rank-\allowbreak{}margin}; \texttt{ass:\allowbreak{}pencil-\allowbreak{}gap}; \texttt{ass:\allowbreak{}source-\allowbreak{}moment-\allowbreak{}envelope}; \texttt{ass:\allowbreak{}mixing-\allowbreak{}singular-\allowbreak{}margin}; \texttt{ass:\allowbreak{}loading-\allowbreak{}selection-\allowbreak{}margin}; \texttt{ass:\allowbreak{}deletion-\allowbreak{}rank-\allowbreak{}margin}.
\end{definition}

\begin{definition}[def:tensor-pencil-inverse]\label{def:tensor-pencil-inverse}
\par\smallskip\noindent\emph{Defined object.} \texttt{Common-\allowbreak{}order tensor-\allowbreak{}pencil inverse}
\par\smallskip\noindent\emph{Construction.} For each \(w\in\{u,v,e_1,\ldots,e_p\}\), contract the last \(q\) modes of \(\mathcal T_r\) against \(u^{\otimes(q-1)}\otimes w\) and matricize the remaining \(2d\) modes as \(M_w=V_d\operatorname{diag}(\lambda_j(u^\top c_j)^{q-1}(w^\top c_j))V_d^\top\). Let \(U\) be an orthonormal basis for \(\operatorname{range}(M_u)\), set \(A_w=U^\top M_wU\) and \(G_w=A_wA_u^{-1}\), form the spectral projector \(P_j\) of \(G_v\), and set \(z_{ij}=\operatorname{tr}(G_{e_i}P_j)\). Return \(c_j=z_j/\lVert z_j\rVert_2\), oriented by \(u^\top c_j>0\), up to a common permutation
\par\smallskip\noindent\emph{Inputs.} \(\mathcal T_r\); \(u\); \(v\).
\end{definition}

\begin{definition}[def:inverse-constants]\label{def:inverse-constants}
\par\smallskip\noindent\emph{Defined object.} \texttt{Explicit local inverse constants}
\par\smallskip\noindent\emph{Construction.} \(\eta=\kappa\sigma^{q+2}\), \(\chi=\sqrt n/\sigma\), \(h=2/\eta+2n\Lambda/\eta^2\), \(\epsilon_0=\min\{\eta/2,\sigma/(6\chi h)\}\), \(L_z=nh\{2\chi+6n\Lambda\chi^2/(\eta\sigma)\}\), and \(L_C=2\sqrt{np}L_z\)
\par\smallskip\noindent\emph{Inputs.} \(p\); \(n\); \(q\); \(\sigma\); \(\kappa\); \(\Lambda\).
\end{definition}

\begin{definition}[def:empirical-cumulant]\label{def:empirical-cumulant}
\par\smallskip\noindent\emph{Defined object.} \texttt{Clipped empirical cumulant tensor}
\par\smallskip\noindent\emph{Construction.} Compute every nonconstant raw empirical moment through total degree \(r\) from \(\mathbf X_N\), clip each coordinate to \([-B_m,B_m]\), and substitute the clipped vector into the moment-cumulant partition polynomial to obtain \(\widehat{\mathcal T}_r\)
\par\smallskip\noindent\emph{Inputs.} \(\mathbf X_N\); \(B_m\).
\end{definition}

\begin{definition}[def:minimum-distance-fit]\label{def:minimum-distance-fit}
\par\smallskip\noindent\emph{Defined object.} \texttt{Separated tensor minimum-\allowbreak{}distance fit}
\par\smallskip\noindent\emph{Construction.} Let \((\widehat C,\widehat\lambda)\) be the lexicographically selected minimizer of \(\lVert\widehat{\mathcal T}_r-\sum_{j=1}^n\lambda'_j(c'_j)^{\otimes r}\rVert_F\) over normalized pairs \((C',\lambda')\) satisfying the numerical restrictions defining \(\mathfrak P_{p,n,r}(a,\delta,\sigma,\kappa,\Lambda,M)\)
\par\smallskip\noindent\emph{Inputs.} \(\widehat{\mathcal T}_r\).
\end{definition}

\begin{definition}[def:estimated-index-set]\label{def:estimated-index-set}
\par\smallskip\noindent\emph{Defined object.} \texttt{Estimated deletion-\allowbreak{}rank set}
\par\smallskip\noindent\emph{Construction.} \(\widehat J_x=\{j\in\mathsf S:\widehat C_{xj}\ne0\text{ and }\operatorname{rank}(\widehat C_{-x,-j})=p-1\}\)
\par\smallskip\noindent\emph{Inputs.} \(\widehat C\); \(x\).
\end{definition}

\begin{definition}[def:confidence-union]\label{def:confidence-union}
\par\smallskip\noindent\emph{Defined object.} \texttt{Finite-\allowbreak{}sample simultaneous confidence union}
\par\smallskip\noindent\emph{Construction.} If \(2\tau_N<\epsilon_0\) and \(b_N<\min\{a,\delta\}/2\), set \(\mathcal U_{xy}=\bigcup_{j\in\widehat J_x}[\widehat C_{yj}/\widehat C_{xj}-w_N,\widehat C_{yj}/\widehat C_{xj}+w_N]\); otherwise set \(\mathcal U_{xy}=[-a^{-1},a^{-1}]\)
\par\smallskip\noindent\emph{Inputs.} \(\widehat C\); \(\widehat J_x\); \(b_N\).
\end{definition}

\begin{definition}[def:multi-order-handle]\label{def:multi-order-handle}
\par\smallskip\noindent\emph{Defined object.} \texttt{Coupled heterogeneous-\allowbreak{}order recovery handle}
\par\smallskip\noindent\emph{Construction.} Given \(P_X\), the known source count \(n\), and an integer cap \(K\ge2n-1\), define \(\mathcal H_{\mathrm{multi}}(P_X,n;K)\) as follows. For every \(\nu\in\{2n-1,\ldots,K\}\), compute \(\mathcal T_\nu(P_X)=\operatorname{cum}_\nu(X)\), set \(\ell_{1,\nu}=\ell_{2,\nu}=n-1\) and \(\ell_{3,\nu}=\nu-2n+2\), and reshape the tensor according to these three blocks as \(\widetilde{\mathcal T}_\nu\). Solve the explicit rank minimization
\[
 R_\nu^\star=
 \min\left\{
 R\in\{0,\ldots,n\}:
 \widetilde{\mathcal T}_\nu
 =\sum_{h=1}^{R}\gamma_h
 \operatorname{vec}(z_h^{\otimes\ell_{1,\nu}})\otimes
 \operatorname{vec}(z_h^{\otimes\ell_{2,\nu}})\otimes
 \operatorname{vec}(z_h^{\otimes\ell_{3,\nu}}),\ 
 \gamma_h\in\mathbb R\setminus\{0\},\ 
 z_h\in\mathbb R^p\setminus\{0\}
 \right\},
\]
where \(R=0\) denotes the empty sum. Retain the projective classes \([z_h]\) from any minimizer, union them over the processed orders using exact projective equality, and return once this union has cardinality \(n\); if it has not reached cardinality \(n\) by \(K\), increase the cap.
\par\smallskip\noindent\emph{Inputs.} \(P_X\); the known source count \(n\); an integer cap \(K\ge2n-1\).
\end{definition}

\begin{definition}[def:observational-law-completion]\label{def:observational-law-completion}
\par\smallskip\noindent\emph{Defined object.} \texttt{Observational-\allowbreak{}law structural completion}
\par\smallskip\noindent\emph{Construction.} For an integer \(m_{\widetilde{\mathcal M}}\ge p\), an observational-law structural completion is a tuple \(\widetilde{\mathcal M}=(m_{\widetilde{\mathcal M}},Q^{\widetilde{\mathcal M}},H^{\widetilde{\mathcal M}},\xi^{\widetilde{\mathcal M}})\), where \(Q^{\widetilde{\mathcal M}},H^{\widetilde{\mathcal M}}\in\mathbb R^{m_{\widetilde{\mathcal M}}\times m_{\widetilde{\mathcal M}}}\), \(\xi^{\widetilde{\mathcal M}}\in\mathbb R^{m_{\widetilde{\mathcal M}}}\), and the endogenous equilibrium satisfies \(Q^{\widetilde{\mathcal M}}V^{\widetilde{\mathcal M}}=H^{\widetilde{\mathcal M}}\xi^{\widetilde{\mathcal M}}\). Its observed mixing matrix is
\[
 A^{\widetilde{\mathcal M}}
 =\bigl[(Q^{\widetilde{\mathcal M}})^{-1}H^{\widetilde{\mathcal M}}\bigr]_{\mathsf O,:}
 \in\mathbb R^{p\times m_{\widetilde{\mathcal M}}}.
\]
\par\smallskip\noindent\emph{Inputs.} \(m_{\widetilde{\mathcal M}}\); \(Q^{\widetilde{\mathcal M}}\); \(H^{\widetilde{\mathcal M}}\); \(\xi^{\widetilde{\mathcal M}}\).
\end{definition}

\begin{definition}[def:law-completion-fiber]\label{def:law-completion-fiber}
\par\smallskip\noindent\emph{Defined object.} \texttt{Source-\allowbreak{}minimal observational-\allowbreak{}law completion fiber}
\par\smallskip\noindent\emph{Construction.} For \(P_X=\mathcal L(C\varepsilon)\), let \(\mathfrak F_x^{\mathrm{law}}(P_X)\) be the set of all observational-law completions \(\widetilde{\mathcal M}\) such that: \(Q^{\widetilde{\mathcal M}}_{ii}=1\) for \(1\le i\le m_{\widetilde{\mathcal M}}\); \(H^{\widetilde{\mathcal M}}\) is monomial; \(\det(Q^{\widetilde{\mathcal M}})\ne0\); \(\det(Q^{\widetilde{\mathcal M}}_{-x,-x})\ne0\); the coordinates of \(\xi^{\widetilde{\mathcal M}}\) are mutually independent, nondegenerate, and non-Gaussian; \(A^{\widetilde{\mathcal M}}\) has no zero column and no two proportional columns; and
\[
 \mathcal L\bigl(A^{\widetilde{\mathcal M}}\xi^{\widetilde{\mathcal M}}\bigr)=P_X.
\]
No spectral-radius or dynamical-stability restriction is imposed.
\par\smallskip\noindent\emph{Member properties.} .
\end{definition}

\begin{definition}[def:law-assigned-source]\label{def:law-assigned-source}
\par\smallskip\noindent\emph{Defined object.} \texttt{Assigned source in a law completion}
\par\smallskip\noindent\emph{Construction.} For \(\widetilde{\mathcal M}\in\mathfrak F_x^{\mathrm{law}}(P_X)\) and \(i\in\mathsf O\), define \(s_{\widetilde{\mathcal M}}^{\mathrm{law}}(i)\) to be the unique \(k\in\{1,\ldots,m_{\widetilde{\mathcal M}}\}\) such that \(H^{\widetilde{\mathcal M}}_{ik}\ne0\).
\par\smallskip\noindent\emph{Inputs.} \(\widetilde{\mathcal M}\); \(i\).
\end{definition}

\begin{definition}[def:law-aligned-source]\label{def:law-aligned-source}
\par\smallskip\noindent\emph{Defined object.} \texttt{Reference-\allowbreak{}aligned source in a law completion}
\par\smallskip\noindent\emph{Construction.} For \(\widetilde{\mathcal M}\in\mathfrak F_x^{\mathrm{law}}(P_X)\) and \(i\in\mathsf O\), define the aligned source \(\bar s_{\widetilde{\mathcal M}}(i)\) to be the unique \(j\in\mathsf S\) such that
\[
 A^{\widetilde{\mathcal M}}_{:,s_{\widetilde{\mathcal M}}^{\mathrm{law}}(i)}
 \in\operatorname{span}(c_j).
\]
\par\smallskip\noindent\emph{Inputs.} \(\widetilde{\mathcal M}\); \(i\); \(C\).
\end{definition}

\begin{definition}[def:observed-response-vector]\label{def:observed-response-vector}
\par\smallskip\noindent\emph{Defined object.} \texttt{Observed hard-\allowbreak{}intervention response vector}
\par\smallskip\noindent\emph{Construction.} For \(\widetilde{\mathcal M}\in\mathfrak F_x^{\mathrm{law}}(P_X)\), define its observed hard-intervention equilibrium response vector by
\[
 \rho_x(\widetilde{\mathcal M})
 =\left(
 \frac{\partial}{\partial t}\operatorname{Eq}_{\widetilde{\mathcal M}}
 \bigl(V_i^{\widetilde{\mathcal M}}\mid\operatorname{do}(V_x^{\widetilde{\mathcal M}}=t)\bigr)
 \right)_{i\in\mathsf O\setminus\{x\}}
 \in\mathbb R^{p-1}.
\]
\par\smallskip\noindent\emph{Inputs.} \(\widetilde{\mathcal M}\); \(x\).
\end{definition}

\begin{definition}[def:response-vector-frontier]\label{def:response-vector-frontier}
\par\smallskip\noindent\emph{Defined object.} \texttt{Deletion-\allowbreak{}rank response-\allowbreak{}vector frontier}
\par\smallskip\noindent\emph{Construction.} Define the response-vector frontier
\[
 \mathcal V_x(C)
 =\left\{
 \left(\frac{C_{ij}}{C_{xj}}\right)_{i\in\mathsf O\setminus\{x\}}
 :j\in J_x(C)
 \right\}\subseteq\mathbb R^{p-1}.
\]
\par\smallskip\noindent\emph{Inputs.} \(C\); \(x\).
\end{definition}

\begin{definition}[def:vector-confidence-union]\label{def:vector-confidence-union}
\par\smallskip\noindent\emph{Defined object.} \texttt{Finite-\allowbreak{}sample response-\allowbreak{}vector confidence union}
\par\smallskip\noindent\emph{Construction.} With \(w_N^{\mathrm{vec}}=\sqrt{p-1}\,w_N\), let \(\overline B_2(z,t)=\{v\in\mathbb R^{p-1}:\lVert v-z\rVert_2\le t\}\). If \(2\tau_N<\epsilon_0\) and \(b_N<\min\{a,\delta\}/2\), define
\[
 \mathcal U_x^{\mathrm{vec}}
 =\bigcup_{j\in\widehat J_x}
 \overline B_2\!\left(
 \left(\frac{\widehat C_{ij}}{\widehat C_{xj}}\right)_{i\in\mathsf O\setminus\{x\}},
 w_N^{\mathrm{vec}}
 \right).
\]
Otherwise define
\[
 \mathcal U_x^{\mathrm{vec}}
 =\overline B_2\!\left(0,\frac{\sqrt{p-1}}a\right).
\]
\par\smallskip\noindent\emph{Inputs.} \(\widehat C\); \(\widehat J_x\); \(b_N\).
\end{definition}

\begin{definition}[def:law-projective-direction-functional]\label{def:law-projective-direction-functional}
\par\smallskip\noindent\emph{Defined object.} \texttt{Observational-\allowbreak{}law projective direction functional}
\par\smallskip\noindent\emph{Construction.} Let \(\mathbb P^{p-1}\) denote real projective space. Define
\[
 \mathcal D(P_X)
 =\bigcup\left\{
   \{[a_k]:1\le k\le m\}:
   \begin{array}{l}
   m\ge p,\ A=[a_1\;\cdots\;a_m]\in\mathbb R^{p\times m},\ \\
   a_k\ne0\text{ and }a_k\notin\operatorname{span}(a_\ell)\text{ for }k\ne\ell,\\
   \xi_1,\ldots,\xi_m\text{ are mutually independent, nondegenerate, and non-Gaussian},\\
   \mathcal L(A\xi)=P_X
   \end{array}
 \right\}.
\]
Thus \(\mathcal D(P_X)\) is defined set-theoretically from the observational law by unioning the unordered projective column sets of all irreducible, all-non-Gaussian independent-component representations. Its asserted uniqueness and equality to the directions of a particular representation are theorem conclusions, not parts of this definition, and the definition does not prescribe an evaluation algorithm.
\par\smallskip\noindent\emph{Inputs.} \(P_X\).
\end{definition}

\section{Main results}

\begin{lemma}[lem:intervention-ratio]\label{lem:intervention-ratio}
For every \(\mathcal M=(Q,H)\in\mathfrak F_x(C,\varepsilon)\), if \(j=s_{\mathcal M}(x)\), then \(C_{xj}\ne0\) and \(\theta_{y\leftarrow x}(\mathcal M)=C_{yj}/C_{xj}\)
\end{lemma}
\begin{proof}
Fix \(\mathcal M=(Q,H)\in\mathfrak F_x(C,\varepsilon)\), put \(j=s_{\mathcal M}(x)\), and write \(h=H_{xj}\). By the monomial-source condition in \(\mathfrak F_x(C,\varepsilon)\), \(h\ne0\) and column \(j\) of \(H\) has no other nonzero entry. The fixed-law identity therefore gives, for every observed \(i\),
\[
 C_{ij}=(Q^{-1}H)_{ij}=(Q^{-1})_{ix}h. \tag{1}
\]
Observational solvability gives \(\det(Q)\ne0\), and postintervention solvability gives \(\det(Q_{-x,-x})\ne0\). The adjugate identity, applied to the \((x,x)\) entry, yields
\[
 (Q^{-1})_{xx}=\frac{\det(Q_{-x,-x})}{\det(Q)}\ne0. \tag{2}
\]
Taking \(i=x\) in (1) and using (2) proves \(C_{xj}\ne0\). Taking \(i=y\) in (1), dividing the two identities by the nonzero quantity \((Q^{-1})_{xx}h\), and then invoking the definition of the equilibrium effect gives
\[
 \frac{C_{yj}}{C_{xj}}
 =\frac{(Q^{-1})_{yx}h}{(Q^{-1})_{xx}h}
 =\frac{(Q^{-1})_{yx}}{(Q^{-1})_{xx}}
 =\theta_{y\leftarrow x}(\mathcal M).
\]
Every division is justified by (2) and \(h\ne0\).
\end{proof}
\par\smallskip\noindent \textit{Justification.} This identity turns a cyclic intervention query into a coordinate of one observed mixing direction. \textit{Gap.} LacerdaSpirtesRamseyHoyer2008 derives square-ICA compatible systems, while no located result states this ratio over an arbitrary-latent fixed-law completion fiber. \textit{Consumer.} It permits an analyst to compute each candidate equilibrium effect without estimating or enumerating a structural graph.

\begin{lemma}[lem:necessary-deletion-rank]\label{lem:necessary-deletion-rank}
For every \(\mathcal M\in\mathfrak F_x(C,\varepsilon)\), its assigned source satisfies \(s_{\mathcal M}(x)\in J_x(C)\)
\end{lemma}
\begin{proof}
Fix \(\mathcal M=(Q,H)\in\mathfrak F_x(C,\varepsilon)\), set \(j=s_{\mathcal M}(x)\), and let \(h=H_{xj}\ne0\). Define
\[
 T=Q^{-1}H,\qquad W=T^{-1}=H^{-1}Q.
\]
By the fixed-law condition, \(T_{\mathsf O,:}=C\). Since \(H\) is monomial and its entry in row \(x\), column \(j\) is \(h\), row \(j\) of \(H^{-1}\) has its unique nonzero entry \(h^{-1}\) in column \(x\). The unit structural diagonal consequently implies
\[
 W_{jx}=(H^{-1}Q)_{jx}=h^{-1}Q_{xx}=h^{-1}\ne0. \tag{3}
\]
The inverse-cofactor formula for \(W=T^{-1}\) now gives
\[
 W_{jx}=\frac{(-1)^{j+x}\det(T_{-x,-j})}{\det(T)}.
\]
Because both \(W_{jx}\) and \(\det(T)\) are nonzero, \(T_{-x,-j}\) is nonsingular. Its rows indexed by \(\mathsf O\setminus\{x\}\) are precisely the \(p-1\) rows of \(C_{-x,-j}\); a subset of the rows of a nonsingular matrix is linearly independent. Hence
\[
 \operatorname{rank}(C_{-x,-j})=p-1. \tag{4}
\]
Finally, as in the adjugate calculation for the intervention ratio,
\[
 C_{xj}=(Q^{-1})_{xx}h
 =\frac{\det(Q_{-x,-x})}{\det(Q)}h\ne0,
\]
using observational and postintervention solvability. Equations (4) and the last display are exactly the two membership conditions for \(j\in J_x(C)\).
\end{proof}
\par\smallskip\noindent \textit{Justification.} Deletion rank is the observable obstruction that rules out source assignments incompatible with a legal intervention. \textit{Gap.} DaiAlbrechtSpirtesZhang2026 uses graph-level edge ranks and children bases, not numerical row-and-column deletion ranks for a fixed law. \textit{Consumer.} It provides a falsifiable certificate that removes impossible effect values before any completion is constructed.

\begin{lemma}[lem:shear-matching-sufficiency]\label{lem:shear-matching-sufficiency}
For every \(j\in J_x(C)\), \(\operatorname{Cert}(C,x)\) returns a completion \(\mathcal M_j\in\mathfrak F_x(C,\varepsilon)\) with \(s_{\mathcal M_j}(x)=j\). If \(C\) is rational, the returned \(Q_j\) and \(H_j\) are rational.
\end{lemma}
\begin{proof}
By \(\textnormal{ass:full-row-rank}\), the \(p\) rows of \(C\) are linearly independent.  Extend them to a basis of \(\mathbb R^n\), as prescribed by \(\textnormal{def:witness-algorithm}\), to obtain \(T_0\in\operatorname{GL}_n(\mathbb R)\) with \((T_0)_{\mathsf O,:}=C\), and put \(W_0=T_0^{-1}\).

We first verify the deletion test used by the algorithm.  Let \(R=C_{-x,:}\).  Its \(p-1\) rows are linearly independent.  For every \(j\in\mathsf S\),
\[
 \operatorname{rank}(R_{:,-j})<p-1
 \quad\Longleftrightarrow\quad
 e_j^\top\in\operatorname{rowspan}(R).                                      \tag{1}
\]
Indeed, rank loss after deleting coordinate \(j\) means that some nonzero linear combination of the rows of \(R\) vanishes outside coordinate \(j\).  That combination is nonzero before deletion, because the rows of \(R\) are independent, and hence it is a nonzero multiple of \(e_j^\top\).  This proves the forward implication in (1), and the reverse implication follows by deleting coordinate \(j\).

The identity \(W_0T_0=I_n\) gives the unique basis expansion
\[
 e_j^\top=\sum_{i=1}^n (W_0)_{ji}(T_0)_{i,:}.                              \tag{2}
\]
The right side of (2) belongs to the span of the observed rows other than \(x\) exactly when its coefficients on row \(x\) and on all rows in \(\mathsf L\) vanish.  Combining this observation with (1) yields
\[
 \operatorname{rank}(C_{-x,-j})=p-1
 \quad\Longleftrightarrow\quad
 \bigl((W_0)_{jx},(W_0)_{j,\mathsf L}\bigr)\ne0.                            \tag{3}
\]
Moreover, \(T_0W_0=I_n\) shows that each column of
\[
 K_x=(W_0)_{:,\{x\}\cup\mathsf L}
\]
lies in \(\ker(C_{-x,:})\).  These \(n-p+1\) columns are independent, while rank--nullity and \(\operatorname{rank}(C_{-x,:})=p-1\) give
\(\dim\ker(C_{-x,:})=n-p+1\).  Thus they form a basis.  In particular, (3) proves the row-nonvanishing test for \(J_x(C)\) in \(\textnormal{def:witness-algorithm}\).

We next justify the single shared matching and its alternating forest.  Since \(W_0\) is invertible, \((W_0)_{:,-x}\) has column rank \(n-1\).  Hence some \((n-1)\)-row minor is nonsingular.  A nonzero term in the determinant expansion of that minor gives a matching \(\mathcal P_0\) in the nonzero-pattern graph of \((W_0)_{:,-x}\) that saturates all columns and leaves one row, say \(j_0\), unmatched.

For any \(j\in J_x(C)\), \(\textnormal{def:admissible-source-set}\) gives \(C_{xj}\ne0\).  The adjugate identity for \(W_0\), together with \(W_0^{-1}=T_0\) and \((T_0)_{xj}=C_{xj}\), gives
\[
 (-1)^{j+x}\det((W_0)_{-j,-x})
 =\det(W_0)(W_0^{-1})_{xj}
 =\det(W_0)C_{xj}\ne0.                                                      \tag{4}
\]
The determinant expansion in (4) therefore supplies another column-saturating matching \(\mathcal P_j\), now leaving row \(j\) unmatched.  In the symmetric difference \(\mathcal P_0\mathbin{\triangle}\mathcal P_j\), every column has degree zero or two, every row other than \(j_0,j\) has degree zero or two, row \(j_0\) has degree one unless \(j=j_0\), and row \(j\) has degree one unless \(j=j_0\).  Consequently its noncyclic component is an alternating path from \(j_0\) to \(j\), with the empty path allowed when \(j=j_0\).  Thus \(j\) is in the alternating-reachability forest rooted at \(j_0\), and flipping \(\mathcal P_0\) along this path produces a matching that saturates every column other than \(x\) and leaves exactly row \(j\) unmatched.  This proves that the one matching and one forest specified in \(\textnormal{def:witness-algorithm}\) suffice simultaneously for all \(j\in J_x(C)\).

Fix such a \(j\).  If \((W_0)_{jx}\ne0\), set \(T_j=T_0\) and \(W_j=W_0\).  Otherwise, (3) provides an \(\ell\in\mathsf L\) with \((W_0)_{j\ell}\ne0\).  Set
\[
 T_j=(I_n+e_\ell e_x^\top)T_0,
 \qquad
 W_j=W_0(I_n-e_\ell e_x^\top).                                             \tag{5}
\]
Because \(x\ne\ell\), one has \((e_\ell e_x^\top)^2=0\), so the two shear factors in (5) are mutual inverses and \(W_j=T_j^{-1}\).  Left multiplication in the definition of \(T_j\) changes only latent row \(\ell\), whence
\[
 (T_j)_{\mathsf O,:}=C.                                                     \tag{6}
\]
Right multiplication in the definition of \(W_j\) changes only column \(x\), and hence
\[
 (W_j)_{:,-x}=(W_0)_{:,-x},
 \qquad
 (W_j)_{jx}=-(W_0)_{j\ell}\ne0.                                            \tag{7}
\]
Thus the matching obtained from the shared alternating forest remains a matching in the nonzero pattern of \((W_j)_{:,-x}\).  If \(n=p\), then \(\mathsf L=\varnothing\), and (3) shows that the second branch is impossible.

Since \(W_j^{-1}=T_j\), equations (6) and \(\textnormal{def:admissible-source-set}\) give another adjugate identity,
\[
 (-1)^{j+x}\det((W_j)_{-j,-x})
 =\det(W_j)(W_j^{-1})_{xj}
 =\det(W_j)C_{xj}\ne0.                                                      \tag{8}
\]
Let \(\mu_j\) assign to each column \(i\ne x\) the row matched to it after the path flip, and extend it by \(\mu_j(x)=j\).  This is a bijection of \(\mathsf S\).  Let \(\Pi_j\) be the permutation matrix whose row \(i\) is \(e_{\mu_j(i)}^\top\).  Every diagonal entry of \(\Pi_jW_j\) is nonzero.  Therefore
\[
 D_j=\operatorname{diag}(\Pi_jW_j),\qquad
 Q_j=D_j^{-1}\Pi_jW_j,\qquad
 H_j=D_j^{-1}\Pi_j                                                       \tag{9}
\]
are well defined.  Equation (9) makes \(H_j\) monomial, makes every diagonal entry of \(Q_j\) equal to one, and makes \(Q_j\) invertible.  Direct multiplication gives
\[
 Q_j^{-1}H_j
 =W_j^{-1}\Pi_j^{-1}D_jD_j^{-1}\Pi_j
 =T_j,
\]
so (6) proves
\[
 [Q_j^{-1}H_j]_{\mathsf O,:}=C.                                            \tag{10}
\]
Because row \(x\) of \(\Pi_j\), and hence row \(x\) of \(H_j\), selects source row \(j\), \(\textnormal{def:assigned-source}\) gives \(s_{\mathcal M_j}(x)=j\).

It remains to verify postintervention solvability rather than infer it from observational invertibility.  After deleting row and column \(x\), \((\Pi_jW_j)_{-x,-x}\) is a row permutation of \((W_j)_{-j,-x}\), because the matching uses all source rows except \(j\) on the columns other than \(x\).  By (8) it is nonsingular.  Since \(D_{j,-x,-x}\) is also nonsingular, (9) yields
\[
 \det((Q_j)_{-x,-x})\ne0.                                                   \tag{11}
\]
Equations (9)--(11) verify, one condition at a time, every requirement of \(\textnormal{def:completion-fiber}\).  Hence \(\mathcal M_j=(Q_j,H_j)\in\mathfrak F_x(C,\varepsilon)\), with the claimed assigned source.

Finally suppose \(C\) is rational.  Because the rational coordinate rows span \(\mathbb R^n\), Gaussian elimination can append \(n-p\) coordinate rows to the independent rational rows of \(C\) to produce a rational \(T_0\).  Its inverse \(W_0\), the shear matrices in (5), and the permutations are rational.  Every diagonal entry inverted in (9) is a nonzero rational number.  Therefore \(Q_j\) and \(H_j\) are rational, as claimed.
\end{proof}
\par\smallskip\noindent \textit{Justification.} This is the nontrivial converse: one latent-row shear and one perfect matching realize every admissible ratio without changing the observed law. \textit{Gap.} The matching construction of DaiAlbrechtSpirtesZhang2026 concerns compatible graph families; LacerdaSpirtesRamseyHoyer2008 has square row permutations and no latent-row shear. \textit{Consumer.} It yields exact audit models demonstrating that an apparent effect ambiguity is structural rather than numerical.

\begin{theorem}[thm:sharp-effect-frontier]\label{thm:sharp-effect-frontier}
Let \(P_X=\mathcal L(C\varepsilon)\). Under the population full-row-rank, nonzero distinct-direction, and independent nondegenerate non-Gaussian source assumptions, the set-valued functional \(\mathcal D(P_X)\) is the unique unordered projective column set common to every irreducible, all-non-Gaussian independent-component representation of \(P_X\), and
\[
 \mathcal D(P_X)=\{[c_j]:j\in\mathsf S\}.
\]
This is a full-law identification statement, not an evaluation algorithm. Moreover, \(J_x(C)\ne\varnothing\), every \(\widetilde{\mathcal M}\in\mathfrak F_x^{\mathrm{law}}(P_X)\) has \(m_{\widetilde{\mathcal M}}=n\), and the sharp law-fiber response-vector set is
\[
 \left\{\rho_x(\widetilde{\mathcal M}):\widetilde{\mathcal M}\in\mathfrak F_x^{\mathrm{law}}(P_X)\right\}
 =\mathcal V_x(C)
 =\left\{\left(\frac{C_{ij}}{C_{xj}}\right)_{i\in\mathsf O\setminus\{x\}}:C_{xj}\ne0,\ \operatorname{rank}(C_{-x,-j})=p-1\right\}.
\]
For every ordered \(y\ne x\), its scalar coordinate projection is exactly
\[
 \left\{[\rho_x(\widetilde{\mathcal M})]_y:\widetilde{\mathcal M}\in\mathfrak F_x^{\mathrm{law}}(P_X)\right\}
 =\mathcal R_{xy}(C).
\]
The frontiers are independent of the representative of \(\mathcal D(P_X)\), because column permutation and nonzero rescaling leave them invariant. The entire response vector is identified if and only if \(\mathcal V_x(C)\) is a singleton, and the \(y\)-coordinate effect is identified if and only if \(\mathcal R_{xy}(C)\) is a singleton. If the latter is not a singleton, two indices with distinct ratios yield two exact same-law SCM countermodels via the canonical embeddings \(\widetilde{\mathcal M}_j=(n,Q_j,H_j,\varepsilon)\) of the shared-kernel and shared-matching certificate construction. Conditional on any exact irreducible mixing representative \(C\), all frontier values and one exact witness per value are computable in \(O(n^3)\) exact real-arithmetic operations. This order is optimal for dense all-witness output in the worst case. This is not an end-to-end complexity claim for obtaining \(C\) from \(P_X\). No stability assumption is imposed.
\end{theorem}
\begin{proof}
The causal target is the set of equilibrium responses over \(\mathfrak F_x^{\mathrm{law}}(P_X)\); the deletion-rank sets are observational-law functionals only after their equality to that target and their invariance to the choice of exact mixing representative have been proved.

We first identify the projective directions from the full observational law. The displayed representation \(P_X=\mathcal L(C\varepsilon)\) is one of the representations unioned in \(\mathcal D(P_X)\): the nonzero-column and distinct-direction assumptions make \(C\) irreducible, while independence, nondegeneracy, and non-Gaussianity give the required source properties. Let \(P_X=\mathcal L(A\xi)\) be any other representation in that union, with \(A\in\mathbb R^{p\times m}\). The hypotheses of \(\mathrm{lem{:}dai\text{-}irreducible\text{-}oica\text{-}uniqueness}\) are exactly the just-listed assumptions, so that lemma yields
\[
 m=n,
 \qquad
 A=C\Pi D, \tag{11}
\]
for a permutation matrix \(\Pi\) and a nonsingular diagonal matrix \(D\). Therefore
\[
 \{[a_k]:1\le k\le m\}=\{[c_j]:j\in\mathsf S\}. \tag{12}
\]
Equation (12) holds for every member of the defining union, and the reference representation supplies one member. By the definition of \(\mathcal D(P_X)\),
\[
 \mathcal D(P_X)=\{[c_j]:j\in\mathsf S\}. \tag{13}
\]
This is uniqueness as a set-valued functional of the full law; no step is an algorithm for evaluating that functional.

We next show that the admissible set is nonempty. Put \(A_x=C_{-x,:}\). Full row rank of \(C\) implies \(\operatorname{rank}(A_x)=p-1\), so \(\dim\ker(A_x)=n-p+1\), whereas \(\dim\ker(C)=n-p\). Hence there exists \(z\in\ker(A_x)\setminus\ker(C)\). Since \(A_xz=0\), the inequality \(Cz\ne0\) can occur only in row \(x\), and thus
\[
 0\ne(Cz)_x=\sum_{j=1}^n C_{xj}z_j. \tag{14}
\]
At least one summand in (14) has \(C_{xj}\ne0\) and \(z_j\ne0\). By the kernel-row characterization in \(\mathrm{lem{:}shared\text{-}kernel\text{-}matching\text{-}certificate}\), this \(j\) belongs to \(J_x(C)\). Therefore \(J_x(C)\ne\varnothing\).

For the necessary inclusion, fix \(\widetilde{\mathcal M}\in\mathfrak F_x^{\mathrm{law}}(P_X)\). The hypotheses of \(\mathrm{lem{:}law\text{-}fiber\text{-}response\text{-}rank}\) are the declared nonzero-column, distinct-direction, independent-source, source-nondegeneracy, and source-non-Gaussianity assumptions. That lemma gives
\[
 m_{\widetilde{\mathcal M}}=n,
 \qquad
 j=\bar s_{\widetilde{\mathcal M}}(x)\in J_x(C),
 \qquad
 \rho_x(\widetilde{\mathcal M})
 =\left(\frac{C_{ij}}{C_{xj}}\right)_{i\in\mathsf O\setminus\{x\}}. \tag{15}
\]
In particular, \(C_{xj}\ne0\), \(\operatorname{rank}(C_{-x,-j})=p-1\), and every response in the law fiber belongs to \(\mathcal V_x(C)\).

For the reverse inclusion, take \(j\in J_x(C)\). The shared-certificate lemma returns \(\mathcal M_j=(Q_j,H_j)\in\mathfrak F_x(C,\varepsilon)\) with \(s_{\mathcal M_j}(x)=j\). Define the observational-law completion
\[
 \widetilde{\mathcal M}_j=(n,Q_j,H_j,\varepsilon). \tag{16}
\]
Membership of \(\mathcal M_j\) in the fixed-law fiber supplies the unit diagonal of \(Q_j\), monomiality of \(H_j\), observational solvability, postintervention solvability, and
\[
 [Q_j^{-1}H_j]_{\mathsf O,:}=C. \tag{17}
\]
The independent, nondegenerate, and non-Gaussian source assumptions apply to \(\varepsilon\), and the nonzero-column and distinct-direction assumptions apply to the observed mixing matrix in (17). Consequently every defining condition of \(\mathfrak F_x^{\mathrm{law}}(P_X)\) holds: in particular,
\[
 \mathcal L([Q_j^{-1}H_j]_{\mathsf O,:}\varepsilon)
 =\mathcal L(C\varepsilon)=P_X. \tag{18}
\]
Thus \(\widetilde{\mathcal M}_j\in\mathfrak F_x^{\mathrm{law}}(P_X)\). Because its exact observed mixing matrix is \(C\), distinct directions make its aligned source at row \(x\) equal to \(j\). Applying \(\mathrm{lem{:}law\text{-}hard\text{-}intervention\text{-}response}\) and then using the unique nonzero entry of row \(x\) of the monomial matrix \(H_j\), equation (17) gives, for every \(i\in\mathsf O\setminus\{x\}\),
\[
 [\rho_x(\widetilde{\mathcal M}_j)]_i
 =\frac{(Q_j^{-1})_{ix}}{(Q_j^{-1})_{xx}}
 =\frac{(Q_j^{-1}H_j)_{ij}}{(Q_j^{-1}H_j)_{xj}}
 =\frac{C_{ij}}{C_{xj}}. \tag{19}
\]
Division is valid because \(j\in J_x(C)\) entails \(C_{xj}\ne0\). Hence every vector in \(\mathcal V_x(C)\) is attained in the law fiber. Together with (15), this proves
\[
 \{\rho_x(\widetilde{\mathcal M}):\widetilde{\mathcal M}\in\mathfrak F_x^{\mathrm{law}}(P_X)\}
 =\mathcal V_x(C). \tag{20}
\]

Coordinate projection at \(y\) sends the vector indexed by \(j\) in (20) to \(C_{yj}/C_{xj}\). The definition of \(\mathcal R_{xy}(C)\) therefore yields
\[
 \{[\rho_x(\widetilde{\mathcal M})]_y:\widetilde{\mathcal M}\in\mathfrak F_x^{\mathrm{law}}(P_X)\}
 =\mathcal R_{xy}(C). \tag{21}
\]
If \(C'\) is another exact representative obtained by selecting one nonzero vector from each member of the identified set (13), then after ordering its columns \(C'=C\Pi D\). The hypotheses of \(\mathrm{lem{:}frontier\text{-}monomial\text{-}invariance}\) hold, and that lemma gives
\[
 \mathcal V_x(C')=\mathcal V_x(C),
 \qquad
 \mathcal R_{xy}(C')=\mathcal R_{xy}(C). \tag{22}
\]
Equations (13) and (22) prove that the rank-and-ratio maps are representative independent.

A response is point identified over a completion fiber exactly when it is constant on that fiber. The sharp equalities (20)--(21) therefore show that the whole response vector is identified if and only if \(\mathcal V_x(C)\) is a singleton, and its \(y\)-coordinate is identified if and only if \(\mathcal R_{xy}(C)\) is a singleton. If the latter set is not a singleton, choose \(j,k\in J_x(C)\) with
\[
 \frac{C_{yj}}{C_{xj}}\ne\frac{C_{yk}}{C_{xk}}. \tag{23}
\]
The two completions (16) for \(j\) and \(k\) both have observational law \(P_X\), both satisfy observational and postintervention solvability by their law-fiber membership, and have distinct \(y\)-coordinate responses by (19) and (23). They are the required exact same-law SCM countermodels.

It remains to certify the complexity assertion. Conditional on exact \(C\), \(\mathrm{lem{:}shared\text{-}kernel\text{-}matching\text{-}certificate}\) computes the single basis extension and inverse, all admissibility decisions from one kernel basis, all leave-one-row matchings from one column-saturating matching and one alternating forest, and materializes every dense \((Q_j,H_j)\) in \(O(n^3)\) exact real-arithmetic operations. Equations (19)--(21) compute the vector and scalar frontier values from the same at most \(n\) columns; exact deduplication remains within \(O(n^3)\). The same lemma certifies that this order is optimal for dense all-witness output in the worst case. Thus the conditional \(O(n^3)\) order is worst-case output optimal. This computation starts from an exact irreducible representative \(C\); equation (13) remains a qualitative full-law identification result and is not an algorithm or complexity bound for recovering \(C\) from \(P_X\). None of the assumptions or constructions imposes spectral-radius or dynamical stability. This completes the proof.
\end{proof}
\par\smallskip\noindent \textit{Justification.} DaiAlbrechtSpirtesZhang2026 maps graphical distributional equivalence in Theorem 2, traversal in Theorem 3, and maximal-class presentation in Theorem 4, but those results do not enumerate the fixed-law postintervention response image, prove deletion-rank equality, or materialize exact same-law effect witnesses. TramontanoKivvaSalehkaleybarDrtonKiyavash2024, Theorem 3.3 gives an if-and-only-if criterion for generic identification of the total causal effect from \(j\) to \(i\) without graph knowledge, but that acyclic selected-effect result does not enumerate the cyclic exact-law response fiber, prove deletion-rank sharpness, or construct same-law countermodels. TramontanoKivvaSalehkaleybarDrtonKiyavash2025, Theorem 3.5 identifies the \(T\)-to-\(Y\) effect generically from the first \(k(\ell)\) cumulants, but that restricted acyclic higher-cumulant result does not characterize all effects across exact same-law cyclic completions, the fixed-law response image, deletion-rank equality, or exact witnesses. TramontanoDrtonEtesami2026 studies coefficient identification conditional on a known mixed graph: its general cyclic result is necessary and its complete result is confined to restricted cycle structures. It neither fixes an irreducible source-direction representation from the observational law nor enumerates the postintervention-solvable same-law completion fiber, so it cannot yield this source-minimal deletion-rank response frontier or its exact witnesses. On the disclosed source-minimal all-non-Gaussian fiber, which is narrower than Dai et al.'s graphically irreducible family, this theorem proves the representative-invariant rank-and-ratio frontier and conditionally constructs every response and dense witness in worst-case output-optimal \(O(n^3)\) exact arithmetic. \textit{Gap.} The precise comparison is theorem-specific. DaiAlbrechtSpirtesZhang2026, Section 4, Theorems 2, 3, and 4 respectively give the graphical equivalence criterion, traversal result, and maximal-class presentation; they do not give the fixed-law postintervention response image, deletion-rank equality, or exact same-law effect witnesses. TramontanoKivvaSalehkaleybarDrtonKiyavash2024, Theorem 3.3 gives the acyclic graph-unknown generic total-effect criterion, but not enumeration of every cyclic exact-law response, deletion-rank sharpness, or same-law countermodels. TramontanoKivvaSalehkaleybarDrtonKiyavash2025, Theorem 3.5 gives generic \(T\)-to-\(Y\) effect identification from the first \(k(\ell)\) cumulants in its restricted acyclic setting, but not all effects across exact same-law cyclic completions, the fixed-law response image, deletion-rank equality, or exact countermodel witnesses. TramontanoDrtonEtesami2026 conditions coefficient identification on a known mixed graph, supplies only a necessary general cyclic criterion and completeness only for restricted cycle structures, and neither uses the source-minimal irreducible OICA law fiber nor enumerates all postintervention-solvable same-law responses. The present theorem fills those latter gaps only for the explicitly source-minimal all-non-Gaussian completion fiber, and does not claim Dai et al.'s larger graphically irreducible scope. \textit{Consumer.} For the fourteen-stock Hong Kong example, this is an exact-\(C\) population certificate with conditional \(O(n^3)\) frontier-and-witness enumeration. Recovering \(C\) empirically, certifying separation, and transferring iid inference to the dependent series require separate evidence.

\begin{proposition}[prop:identification-and-countermodels]\label{prop:identification-and-countermodels}
Let \(P_X=\mathcal L(C\varepsilon)\). Suppose that \(\operatorname{rank}(C)=p\), every column of \(C\) is nonzero, no two columns of \(C\) are proportional, and the coordinates of \(\varepsilon\) are mutually independent, nondegenerate, and non-Gaussian. The law-fiber response vector \(\rho_x\) is point identified over \(\mathfrak F_x^{\mathrm{law}}(P_X)\) if and only if \(|\mathcal V_x(C)|=1\), equivalently \(|J_x(C)|=1\). If this vector criterion fails, there are distinct \(j,k\in J_x(C)\) for which the canonical embeddings of the completions returned by \(\operatorname{Cert}(C,x)\) are two exact same-observational-law members \(\widetilde{\mathcal M}_j,\widetilde{\mathcal M}_k\in\mathfrak F_x^{\mathrm{law}}(P_X)\), both satisfying postintervention solvability, such that \(\rho_x(\widetilde{\mathcal M}_j)\ne\rho_x(\widetilde{\mathcal M}_k)\). They can be constructed in \(O(n^3)\) additional arithmetic operations after the shared basis extension and inverse. For every \(y\in\mathsf O\setminus\{x\}\), the scalar coordinate \([\rho_x]_y\) is point identified over the same law fiber if and only if
\[
 |\mathcal R_{xy}(C)|=1,
 \qquad\text{equivalently}\qquad
 \operatorname{rank}\!\left(C_{\{x,y\},J_x(C)}\right)=1.
\]
If this scalar criterion fails, two canonical witnesses have distinct \(y\)-coordinate responses. Each scalar frontier is the corresponding coordinate projection of \(\mathcal V_x(C)\); the claim does not combine coordinates arising from different source indices into a response vector.
\end{proposition}
\begin{proof}
Define, for \(j\in J_x(C)\),
\[
 \phi_x(j)=
 \left(\frac{C_{ij}}{C_{xj}}\right)_{i\in\mathsf O\setminus\{x\}}.
                                                                            \tag{18}
\]
Every denominator in (18) is nonzero by \(\textnormal{def:admissible-source-set}\).  By \(\textnormal{thm:sharp-effect-frontier}\), the image of the map
\[
 \widetilde{\mathcal M}\longmapsto\rho_x(\widetilde{\mathcal M}),
 \qquad
 \widetilde{\mathcal M}\in\mathfrak F_x^{\mathrm{law}}(P_X),
\]
is the nonempty set
\[
 \mathcal V_x(C)=\{\phi_x(j):j\in J_x(C)\},                                 \tag{19}
\]
and \(J_x(C)\ne\varnothing\).  Therefore the response vector is point identified exactly when \(|\mathcal V_x(C)|=1\).

We show that \(\phi_x\) is injective.  If \(j,k\in J_x(C)\) and
\(\phi_x(j)=\phi_x(k)\), then
\[
 \frac{C_{ij}}{C_{xj}}=\frac{C_{ik}}{C_{xk}}
 \qquad(i\in\mathsf O\setminus\{x\}).                                      \tag{20}
\]
At coordinate \(x\), both normalized columns equal one.  Adjoining that coordinate to (20) gives
\[
 \frac{c_j}{C_{xj}}=\frac{c_k}{C_{xk}},
 \qquad\text{so}\qquad
 c_j=\frac{C_{xj}}{C_{xk}}c_k.                                             \tag{21}
\]
The multiplier is nonzero.  By \(\textnormal{ass:distinct-directions}\), (21) forces \(j=k\).  Hence
\[
 |\mathcal V_x(C)|=|J_x(C)|,                                                \tag{22}
\]
and nonemptiness together with (19)--(22) proves the stated whole-vector criterion.

If this criterion fails, (22) supplies distinct \(j,k\in J_x(C)\) with
\(\phi_x(j)\ne\phi_x(k)\).  Applying
\(\textnormal{lem:canonical-exact-c-law-witness}\) twice gives
\[
 \widetilde{\mathcal M}_j=(n,Q_j,H_j,\varepsilon),\qquad
 \widetilde{\mathcal M}_k=(n,Q_k,H_k,\varepsilon)
 \quad\text{in }\mathfrak F_x^{\mathrm{law}}(P_X)
\]
and
\[
 \rho_x(\widetilde{\mathcal M}_j)=\phi_x(j)
 \ne\phi_x(k)=\rho_x(\widetilde{\mathcal M}_k).                             \tag{23}
\]
Membership in \(\mathfrak F_x^{\mathrm{law}}(P_X)\), by
\(\textnormal{def:law-completion-fiber}\), separately certifies that both systems have observational law \(P_X\), invertible structural matrices, and invertible postintervention minors.  Thus (23) consists of two exact same-law, postintervention-solvable countermodels.

For the scalar claim, \(\textnormal{thm:sharp-effect-frontier}\) and
\(\textnormal{def:effect-frontier}\) give the coordinate projection
\[
 \left\{[\rho_x(\widetilde{\mathcal M})]_y:
 \widetilde{\mathcal M}\in\mathfrak F_x^{\mathrm{law}}(P_X)\right\}
 =\mathcal R_{xy}(C)
 =\left\{\frac{C_{yj}}{C_{xj}}:j\in J_x(C)\right\}.                         \tag{24}
\]
Therefore the \(y\)-coordinate is point identified exactly when
\(|\mathcal R_{xy}(C)|=1\).

Consider \(C_{\{x,y\},J_x(C)}\).  Its \(x\)-row has no zero coordinate, so the matrix has rank at least one.  It has rank at most one exactly when all its \(2\)-by-\(2\) minors vanish:
\[
 C_{xj}C_{yk}-C_{xk}C_{yj}=0
 \qquad(j,k\in J_x(C)).                                                     \tag{25}
\]
Dividing (25) by the nonzero product \(C_{xj}C_{xk}\) shows that (25) is equivalent to
\[
 \frac{C_{yj}}{C_{xj}}=\frac{C_{yk}}{C_{xk}}
 \qquad(j,k\in J_x(C)).                                                     \tag{26}
\]
Because \(J_x(C)\) is nonempty, (24)--(26) prove
\[
 |\mathcal R_{xy}(C)|=1
 \quad\Longleftrightarrow\quad
 \operatorname{rank}\!\left(C_{\{x,y\},J_x(C)}\right)=1.                   \tag{27}
\]
If (27) fails, some \(j,k\in J_x(C)\) violate (26).  Their two canonical witnesses from
\(\textnormal{lem:canonical-exact-c-law-witness}\) then have distinct \(y\)-coordinate responses.  Notice that (19) always assigns a whole response vector using one common source index; (24) is only its coordinate projection and does not form vectors by combining coordinates from different indices.

It remains to certify the stated construction cost.  After the shared basis extension and inverse in
\(\textnormal{def:witness-algorithm}\), forming the nonzero-pattern graph costs \(O(n^2)\) scalar inspections.  An augmenting-path matching implementation performs at most \(n-1\) augmentations, each inspecting \(O(n^2)\) incidences, so a column-saturating matching and its alternating-reachability forest cost \(O(n^3)\) operations.  For each of the two selected indices, a stored path flip costs \(O(n)\), while a possible rank-one shear and the explicit formation of the two dense matrices \(Q_j,H_j\) cost \(O(n^2)\).  Thus the two exact countermodels in (23) require \(O(n^3)\) additional operations after the shared extension and inverse, as asserted.  No stability condition is present in \(\textnormal{def:law-completion-fiber}\) or used above.
\end{proof}
\par\smallskip\noindent \textit{Justification.} Separate vector and scalar singleton tests make the distinction between joint response identification and coordinatewise identification operational. \textit{Gap.} No cited cyclic latent result gives the law-fiber criterion \(|J_x(C)|=1\) for the whole response vector or the two exact same-law witness certificate. \textit{Consumer.} A user receives one coherent response vector or two complete structural models demonstrating ambiguity, while scalar effects remain valid coordinate projections.

\begin{theorem}[thm:square-reduction]\label{thm:square-reduction}
When \(n=p\), the latent-row shear is unnecessary. Section 4 of LacerdaSpirtesRamseyHoyer2008 constructs candidate models from admissible row permutations, and its Theorem 1 states that the output of LiNG-D is a distribution-entailment equivalence class. Among those square admissible row-permutation candidates, restrict to the candidates satisfying \(\det(Q_{-x,-x})\ne0\). The projection of this postintervention-solvable subset onto \(\theta_{y\leftarrow x}\) equals \(\mathcal R_{xy}(C)\).
\end{theorem}
\begin{proof}
Assume \(n=p\).  By \(\textnormal{ass:full-row-rank}\), the square matrix \(C\) is nonsingular.  Since \(\mathsf O=\mathsf S\), the extension in
\(\textnormal{def:witness-algorithm}\) is forced to be
\[
 T_0=C,\qquad W_0=C^{-1},
\]
and \(\mathsf L=\varnothing\).  Hence the latent-row shear cannot occur.

By \(\textnormal{lem:lacerda-square-row-permutation}\), every admissible square row-permutation candidate based on \(W_0\) has
\[
 Q=D^{-1}\Pi W_0,\qquad
 D=\operatorname{diag}(\Pi W_0),                                           \tag{28}
\]
where \(\Pi\) is a row-permutation matrix and every diagonal entry of
\(\Pi W_0\) is nonzero.  Associate with it
\[
 H=D^{-1}\Pi.                                                               \tag{29}
\]
Then \(H\) is monomial, \(Q_{ii}=1\) for every \(i\), and \(Q,H\) are invertible.  Equations (28)--(29) give
\[
 Q^{-1}H
 =W_0^{-1}\Pi^{-1}DD^{-1}\Pi
 =C.                                                                        \tag{30}
\]
Therefore, if the candidate is retained by the stipulated filter
\(\det(Q_{-x,-x})\ne0\), equations (28)--(30) verify every condition of
\(\textnormal{def:completion-fiber}\) for \(\mathcal M=(Q,H)\).

Under \(\textnormal{ass:nonzero-columns}\),
\(\textnormal{ass:distinct-directions}\),
\(\textnormal{ass:independent-sources}\),
\(\textnormal{ass:nondegenerate-sources}\), and
\(\textnormal{ass:nongaussian-sources}\), the canonical embedding
\[
 \widetilde{\mathcal M}=(p,Q,H,\varepsilon)
\]
also belongs to \(\mathfrak F_x^{\mathrm{law}}(P_X)\): equation (30) gives both its exact observed mixing matrix \(C\) and its observed law \(P_X\), while the remaining requirements follow from (28)--(29), the intervention-minor filter, and the five displayed assumptions.

The filter also makes the response denominator nonzero.  Indeed, the adjugate identity gives
\[
 (Q^{-1})_{xx}
 =\frac{\det(Q_{-x,-x})}{\det(Q)}\ne0.                                     \tag{31}
\]
By \(\textnormal{def:equilibrium-effect}\) and
\(\textnormal{lem:law-hard-intervention-response}\),
\[
 \theta_{y\leftarrow x}(\mathcal M)
 =\frac{(Q^{-1})_{yx}}{(Q^{-1})_{xx}}
 =[\rho_x(\widetilde{\mathcal M})]_y.                                      \tag{32}
\]
The law-fiber equality in \(\textnormal{thm:sharp-effect-frontier}\) now implies from (32) that every retained square candidate has
\[
 \theta_{y\leftarrow x}(\mathcal M)\in\mathcal R_{xy}(C).                  \tag{33}
\]
This proves one inclusion.

Conversely, let \(t\in\mathcal R_{xy}(C)\).  By
\(\textnormal{def:effect-frontier}\) and
\(\textnormal{def:admissible-source-set}\), there is \(j\in\mathsf S\) such that
\[
 t=\frac{C_{yj}}{C_{xj}},\qquad
 C_{xj}\ne0,\qquad
 \operatorname{rank}(C_{-x,-j})=p-1.                                      \tag{34}
\]
Because \(n=p\), the last matrix in (34) is square, and its rank condition is equivalent to
\(\det(C_{-x,-j})\ne0\).  Applying the adjugate identity to \(C\) yields
\[
 (W_0)_{jx}
 =(C^{-1})_{jx}
 =\frac{(-1)^{x+j}\det(C_{-x,-j})}{\det(C)}
 \ne0.                                                                      \tag{35}
\]
Thus the certificate takes its no-shear branch for \(j\).  Applying the same identity to \(W_0\) gives
\[
 (-1)^{j+x}\det((W_0)_{-j,-x})
 =\det(W_0)(W_0^{-1})_{xj}
 =\det(W_0)C_{xj}\ne0.                                                      \tag{36}
\]
A nonzero determinant term in (36) gives a perfect matching between the rows other than \(j\) and the columns other than \(x\).  Augment it by assigning row \(j\) to column \(x\), and let \(\Pi_j\) be the resulting row permutation.  Then
\[
 D_j=\operatorname{diag}(\Pi_jW_0),\qquad
 Q_j=D_j^{-1}\Pi_jW_0,\qquad
 H_j=D_j^{-1}\Pi_j                                                        \tag{37}
\]
is exactly a square admissible row-permutation candidate of
\(\textnormal{lem:lacerda-square-row-permutation}\).  After row \(x\) and column \(x\) are deleted, \((\Pi_jW_0)_{-x,-x}\) is a row permutation of the nonsingular matrix in (36); diagonal row rescaling in (37) preserves nonsingularity.  Hence
\[
 \det((Q_j)_{-x,-x})\ne0,                                                   \tag{38}
\]
so this candidate survives the postintervention-solvability filter.  The same membership conclusion is also the \(n=p\) specialization of
\(\textnormal{lem:shear-matching-sufficiency}\).

By \(\textnormal{lem:canonical-exact-c-law-witness}\), its canonical embedding belongs to the law fiber and satisfies
\[
 \rho_x(\widetilde{\mathcal M}_j)
 =\left(\frac{C_{ij}}{C_{xj}}\right)_{i\in\mathsf O\setminus\{x\}}.          \tag{39}
\]
Equations (31)--(32), applied to \(Q_j\), and then (34) and (39), give
\[
 \theta_{y\leftarrow x}(\mathcal M_j)
 =[\rho_x(\widetilde{\mathcal M}_j)]_y
 =\frac{C_{yj}}{C_{xj}}
 =t.                                                                        \tag{40}
\]
Thus every element of \(\mathcal R_{xy}(C)\) is attained by a retained square candidate.  Combining (33) and (40) proves the asserted equality.

Finally, \(\textnormal{thm:sharp-effect-frontier}\) identifies
\(\mathcal R_{xy}(C)\) as the \(y\)-coordinate projection of the complete law-fiber response-vector frontier \(\mathcal V_x(C)\), while
\(\textnormal{lem:lacerda-distribution-entailment-equivalence}\) supplies the distribution-entailment equivalence-class interpretation of the square LiNG-D output.  The argument above does not infer postintervention solvability from observational equivalence: equation (38) verifies it and singular-minor candidates are discarded.  No spectral-radius or dynamical-stability condition was used.
\end{proof}
\par\smallskip\noindent \textit{Justification.} The square reduction shows that the postintervention-solvable LiNG-D row-permutation candidates attain exactly the scalar coordinate projection of the accepted observational-law response-vector frontier. \textit{Gap.} LacerdaSpirtesRamseyHoyer2008 supplies square row-permutation candidates and distribution entailment, but it neither filters a specified intervention minor nor proves that the surviving candidates realize the coordinate projection of the law-fiber response frontier. \textit{Consumer.} Square cyclic-LiNG analyses can enumerate the law-compatible response-vector frontier and answer a scalar query by discarding candidates with singular postintervention systems and projecting the survivors onto the requested outcome coordinate.

\begin{theorem}[thm:uniform-confidence-frontier]\label{thm:uniform-confidence-frontier}
For every \(N\) and every nonempty class \(\mathfrak P_{p,n,r}(a,\delta,\sigma,\kappa,\Lambda,M)\) with \(d\ge1\), \(q\ge1\), and \(n\le\binom{p+d-1}{d}\), the scalar and response-vector confidence unions computed from \(\mathbf X_N\) satisfy
\[
\inf_{\zeta=(P_X,C,\varepsilon,\lambda)\in
\mathfrak P_{p,n,r}(a,\delta,\sigma,\kappa,\Lambda,M)}
P_X^{\otimes N}\!\left\{
\begin{array}{l}
\mathcal V_x(C)\subseteq\mathcal U_x^{\mathrm{vec}}
 \text{ for every }x\in\mathsf O,\quad
\mathcal R_{xy}(C)\subseteq\mathcal U_{xy}
 \text{ for every }x\ne y,\\[2pt]
\text{and, if }2\tau_N<\epsilon_0
 \text{ and }b_N<\min\{a,\delta\}/2,\text{ then after one common permutation},\\[2pt]
\widehat J_x=J_x(C)
 \text{ and }d_H(\mathcal U_x^{\mathrm{vec}},\mathcal V_x(C))
 \le2w_N^{\mathrm{vec}}
 \text{ for every }x\in\mathsf O,\\[2pt]
d_H(\mathcal U_{xy},\mathcal R_{xy}(C))\le2w_N
 \text{ for every }x\ne y
\end{array}
\right\}\ge1-\alpha.
\]
Writing \(\operatorname{pr}_{xy}:\mathbb R^{p-1}\to\mathbb R\) for the \(y\)-coordinate projection associated with treatment \(x\), one has
\[
 \mathcal R_{xy}(C)=\operatorname{pr}_{xy}(\mathcal V_x(C)),
 \qquad
 \mathcal U_{xy}\subseteq
 \operatorname{pr}_{xy}(\mathcal U_x^{\mathrm{vec}}).
\]
Thus the scalar results are coordinatewise reports for the same response coordinates; no Cartesian product of scalar unions is asserted. By the sharp effect-frontier theorem, \(\mathcal V_x(C)\) and \(\mathcal R_{xy}(C)\) are exactly the response-vector and scalar causal targets over \(\mathfrak F_x^{\mathrm{law}}(P_X)\). Uniformly over the same representation-indexed tuples in the informative regime, \(2w_N^{\mathrm{vec}}=O(N^{-1/2})\) and \(2w_N=O(N^{-1/2})\).
\end{theorem}
\begin{proof}
Fix \(\zeta=(P_X,C,\varepsilon,\lambda)\in\mathfrak P_{p,n,r}(a,\delta,\sigma,\kappa,\Lambda,M)\). By the unit-normalization member condition in \(\mathfrak P_{p,n,r}\),
\[
 \lVert X\rVert_2
 =\left\lVert\sum_{j=1}^nc_j\varepsilon_j\right\rVert_2
 \le\sum_{j=1}^n|\varepsilon_j|.
\]
For nonnegative \(t_1,\ldots,t_n\), convexity of \(s\mapsto s^{2r}\) gives
\[
 \left(\sum_{j=1}^nt_j\right)^{2r}
 \le n^{2r-1}\sum_{j=1}^nt_j^{2r}.
\]
The source-moment-envelope assumption and \(M\ge1\) therefore imply
\[
 \mathbb E\lVert X\rVert_2^{2r}
 \le n^{2r-1}\sum_{j=1}^n\mathbb E|\varepsilon_j|^{2r}
 \le n^{2r}M
 =n^{2r}\max\{1,M\}=H.
 \tag{13}
\]

Let \(\mu\in\mathbb R^m\) collect all nonconstant raw moments of \(X\) through total degree \(r\), and let \(\widetilde\mu\) collect their untrimmed empirical averages. For a monomial \(X^\beta\) of degree \(1\le|\beta|\le r\),
\[
 |X^\beta|^2\le1+\lVert X\rVert_2^{2r}.
\]
The iid-sampling assumption and (13) consequently give
\[
 \mathbb E\lVert\widetilde\mu-\mu\rVert_2^2
 =\sum_{k=1}^m\operatorname{Var}(\widetilde\mu_k)
 \le\frac{m(1+H)}N.
 \tag{14}
\]
Cauchy--Schwarz and (13) give \(|\mu_k|\le\sqrt{1+H}\le1+H=B_m\). Coordinatewise clipping to \([-B_m,B_m]\) is Euclidean projection onto an interval containing \(\mu_k\), so the clipped vector \(\widehat\mu\) satisfies
\[
 \lVert\widehat\mu-\mu\rVert_2
 \le\lVert\widetilde\mu-\mu\rVert_2.
\]
For the nonnegative random variable \(Z=\lVert\widehat\mu-\mu\rVert_2^2\), the elementary inequality \(\mathbb EZ\ge t\,P_X^{\otimes N}(Z>t)\), with \(t=e_N^2=m(1+H)/(N\alpha)\), proves
\[
 E_N=\{\lVert\widehat\mu-\mu\rVert_2\le e_N\},
 \qquad
 \inf_{\zeta\in\mathfrak P_{p,n,r}(a,\delta,\sigma,\kappa,\Lambda,M)}
 P_X^{\otimes N}(E_N)\ge1-\alpha.
 \tag{15}
\]

The moment--cumulant partition polynomial used in the definition of \(\widehat{\mathcal T}_r\) is, for each ordered tensor coordinate,
\[
 \operatorname{cum}(X_{i_1},\ldots,X_{i_r})
 =\sum_{\pi\in\mathfrak P_r}(-1)^{|\pi|-1}(|\pi|-1)!
 \prod_{A\in\pi}\mathbb E\!\left[\prod_{k\in A}X_{i_k}\right].
 \tag{16}
\]
For numbers \(a_1,\ldots,a_s,b_1,\ldots,b_s\in[-B_m,B_m]\), telescoping the product gives
\[
 \left|\prod_{h=1}^sa_h-\prod_{h=1}^sb_h\right|
 \le B_m^{s-1}\sum_{h=1}^s|a_h-b_h|
 \le sB_m^{s-1}\lVert\widehat\mu-\mu\rVert_2.
\]
Applying this inequality to every partition in (16), and then summing the squares of the \(p^r\) ordered tensor-coordinate bounds, gives on \(E_N\)
\[
 \begin{aligned}
 \lVert\widehat{\mathcal T}_r-\mathcal T_r\rVert_F
 &\le p^{r/2}
 \sum_{\pi\in\mathfrak P_r}(|\pi|-1)!|\pi|B_m^{|\pi|-1}e_N\\
 &=L_me_N=\tau_N.
 \end{aligned}
 \tag{17}
\]

Let \(K\) be the numerical feasible set in the definition of the separated tensor minimum-distance fit. It is nonempty because the true pair \((C,\lambda)\) is feasible. It is bounded because every column has unit norm and \(\kappa\le|\lambda_j|\le\Lambda\). It is closed because the unit-norm, probe, lifted-singular-value, pencil-gap, mixing-singular-value, loading-alternative, and deletion-singular-value restrictions are inverse images of closed sets under continuous functions. The otherwise nonclosed requirement of pairwise nonproportional columns is redundant on this set: proportional columns have equal ratios \(v^\top c_j/(u^\top c_j)\) and hence violate the positive pencil gap. Thus \(K\) is compact in finite-dimensional Euclidean space. A minimizing sequence for the continuous objective has a convergent subsequence in \(K\), whose limit is a minimizer. Successive minimization of its finitely many coordinates over compact argmin sets supplies the lexicographic selection in the definition.

Define the fitted tensor
\[
 \widehat{\mathcal T}^{\mathrm{fit}}_r
 =\sum_{j=1}^n\widehat\lambda_j\widehat c_j^{\otimes r}.
\]
Since the true pair is feasible, optimality and (17) imply on \(E_N\)
\[
 \begin{aligned}
 \lVert\widehat{\mathcal T}^{\mathrm{fit}}_r-\mathcal T_r\rVert_F
 &\le
 \lVert\widehat{\mathcal T}^{\mathrm{fit}}_r-
 \widehat{\mathcal T}_r\rVert_F
 +\lVert\widehat{\mathcal T}_r-\mathcal T_r\rVert_F\\
 &\le2\lVert\widehat{\mathcal T}_r-\mathcal T_r\rVert_F
 \le2\tau_N.
 \end{aligned}
 \tag{18}
\]
Both \((C,\lambda)\) and \((\widehat C,\widehat\lambda)\) satisfy the numerical hypotheses of the numerical local-inverse lemma. Therefore, whenever \(2\tau_N<\epsilon_0\), equations (18) and that lemma supply one common permutation of the fitted columns under which
\[
 \lVert\widehat C-C\rVert_F
 \le2L_C\tau_N=b_N.
 \tag{19}
\]
All remaining comparisons use this single permutation; in particular, no outcome-specific relabeling is allowed.

Suppose also that \(b_N<\min\{a,\delta\}/2\). Both feasible matrices satisfy the loading alternative that every entry is zero or has magnitude at least \(a\). If corresponding entries of \(C\) and \(\widehat C\) had different zero status, their absolute difference would be at least \(a\), contradicting (19). Their zero patterns therefore agree.

For a matrix \(G\) with \(p-1\) rows, define its least singular value by
\[
 s_{\min}(G)=\inf_{\lVert z\rVert_2=1}\lVert G^\top z\rVert_2.
\]
For any conformable \(G,G'\), the triangle inequality gives
\[
 |s_{\min}(G')-s_{\min}(G)|
 \le\lVert G'-G\rVert_{\mathrm{op}}.
 \tag{20}
\]
Apply (20) to every pair \(\widehat C_{-i,-j},C_{-i,-j}\). Deleting rows and columns cannot increase the Frobenius norm, so (19) bounds the right-hand side by \(b_N<\delta/2\). The deletion-rank-margin assumption, which is imposed on both the true pair and the feasible fit, says that each least singular value is either zero or at least \(\delta\). Thus the two least singular values have the same zero status. Together with equality of the loading zero patterns and the definition of \(J_x\), this proves simultaneously for every \(x\in\mathsf O\) that
\[
 \widehat J_x=J_x(C).
 \tag{21}
\]

Fix \(x\in\mathsf O\) and \(j\in J_x(C)=\widehat J_x\). The loading-selection margin gives \(|C_{xj}|,|\widehat C_{xj}|\ge a\), and unit normalization gives \(|C_{ij}|\le1\). The coordinatewise consequence of (19) therefore yields, for every \(i\ne x\),
\[
 \begin{aligned}
 \left|\frac{\widehat C_{ij}}{\widehat C_{xj}}
 -\frac{C_{ij}}{C_{xj}}\right|
 &\le
 \frac{|\widehat C_{ij}-C_{ij}|}{|\widehat C_{xj}|}
 +\frac{|C_{ij}|\,|\widehat C_{xj}-C_{xj}|}
 {|\widehat C_{xj}|\,|C_{xj}|}\\
 &\le(a^{-1}+a^{-2})b_N=w_N.
 \end{aligned}
 \tag{22}
\]
Set
\[
 r_{xj}=\left(\frac{C_{ij}}{C_{xj}}\right)_{i\in\mathsf O\setminus\{x\}},
 \qquad
 \widehat r_{xj}=\left(
 \frac{\widehat C_{ij}}{\widehat C_{xj}}
 \right)_{i\in\mathsf O\setminus\{x\}}.
\]
Summing the squares of the \(p-1\) bounds in (22) proves
\[
 \lVert\widehat r_{xj}-r_{xj}\rVert_2
 \le\sqrt{p-1}\,w_N=w_N^{\mathrm{vec}}.
 \tag{23}
\]

By (21), every \(r_{xj}\in\mathcal V_x(C)\) has a confidence ball centered at the corresponding \(\widehat r_{xj}\); equation (23) places \(r_{xj}\) in that ball. Hence \(\mathcal V_x(C)\subseteq\mathcal U_x^{\mathrm{vec}}\). Conversely, if \(z\in\mathcal U_x^{\mathrm{vec}}\), some \(j\in J_x(C)\) satisfies \(\lVert z-\widehat r_{xj}\rVert_2\le w_N^{\mathrm{vec}}\), and (23) gives
\[
 \inf_{r\in\mathcal V_x(C)}\lVert z-r\rVert_2
 \le\lVert z-r_{xj}\rVert_2
 \le2w_N^{\mathrm{vec}}.
\]
The two directed-distance bounds therefore prove
\[
 d_H(\mathcal U_x^{\mathrm{vec}},\mathcal V_x(C))
 \le2w_N^{\mathrm{vec}}.
 \tag{24}
\]
The scalar form of (22) gives, by the identical two-directed-distance argument,
\[
 \mathcal R_{xy}(C)\subseteq\mathcal U_{xy},
 \qquad
 d_H(\mathcal U_{xy},\mathcal R_{xy}(C))\le2w_N
 \tag{25}
\]
for every ordered \(x\ne y\).

If either informative-regime inequality fails, the vector definition instead gives the ball of radius \(\sqrt{p-1}/a\) centered at zero. For every \(j\in J_x(C)\), the loading margin, unit normalization, and \(C_{xj}\ne0\) give
\[
 \left\lVert
 \left(\frac{C_{ij}}{C_{xj}}\right)_{i\ne x}
 \right\rVert_2
 \le\frac{\sqrt{p-1}}a.
\]
Thus the fallback ball covers \(\mathcal V_x(C)\). The same argument in one coordinate gives \(|C_{yj}/C_{xj}|\le a^{-1}\), so the scalar fallback interval \([-a^{-1},a^{-1}]\) covers \(\mathcal R_{xy}(C)\).

These conclusions hold simultaneously for all \(x\) and \(y\) on the single event \(E_N\). Equation (15) therefore proves the displayed uniform coverage probability. The quantifier is over the separated representation-indexed iid tuples, and the probability is \(P_X^{\otimes N}\); no sampling uniformity over the qualitative structural completion fiber is used.

For completeness, the response-vector frontier definition gives
\[
 \operatorname{pr}_{xy}(\mathcal V_x(C))=\mathcal R_{xy}(C).
\]
In the informative regime, every point of the scalar interval indexed by \(j\) can be lifted from \(\widehat r_{xj}\) by changing only its \(y\)-coordinate by at most \(w_N\le w_N^{\mathrm{vec}}\); it therefore belongs to the projection of the corresponding Euclidean ball. In the fallback regime,
\[
 [-a^{-1},a^{-1}]
 \subseteq[-\sqrt{p-1}/a,\sqrt{p-1}/a],
\]
the coordinate projection of the fallback ball. Hence \(\mathcal U_{xy}\subseteq\operatorname{pr}_{xy}(\mathcal U_x^{\mathrm{vec}})\). The balls are indexed by one common source \(j\), so the construction never combines coordinates from different completions. The sharp effect-frontier theorem identifies \(\mathcal V_x(C)\) and its scalar projections with the law-fiber causal targets; the confidence sets are observed-sample procedures for those targets, not definitions of them.

Finally, with \(p,n,r,d,q,a,\delta,\sigma,\kappa,\Lambda,M\), and \(\alpha\) fixed, the quantities \(H,B_m,L_m,L_C\) are fixed constants. The definitions give \(e_N=O(N^{-1/2})\), hence \(\tau_N=O(N^{-1/2})\), \(b_N=O(N^{-1/2})\),
\[
 2w_N=O(N^{-1/2}),
 \qquad
 2w_N^{\mathrm{vec}}=2\sqrt{p-1}\,w_N=O(N^{-1/2}),
\]
uniformly over the stated tuple-indexed class.
\end{proof}
\par\smallskip\noindent \textit{Justification.} On the separated common-cumulant-order iid tuple class, the theorem gives simultaneous source-indexed Euclidean balls and scalar intervals with uniform finite-sample coverage and root-\(N\) Hausdorff uncertainty. Its assumptions are the class's fixed-order moment envelope and quantitative separation and causal-map margins, not moment determinacy. \textit{Gap.} Published direction-recovery perturbation bounds do not preserve completion-wise response-vector coupling or certify deletion-rank-selected causal frontiers with simultaneous finite-sample coverage. \textit{Consumer.} Separated iid applications can report uncertainty around each attainable response vector without shrinking the genuine between-completion ambiguity; this theorem does not by itself justify inference for dependent financial time series.

\begin{lemma}[lem:gvx-robust-tensor-scope]\label{lem:gvx-robust-tensor-scope}
\texttt{GoyalVempalaXiao2014 is a robustness anchor for overcomplete tensor decomposition under quantitative nondegeneracy conditions;\allowbreak{} no uniform robustness theorem from that paper is imported into the unrestricted heterogeneous-\allowbreak{}order class here.}
\end{lemma}

\begin{lemma}[lem:lacerda-square-row-permutation]\label{lem:lacerda-square-row-permutation}
\texttt{For a square ICA unmixing matrix,\allowbreak{} Section 4 of LacerdaSpirtesRamseyHoyer2008 constructs candidate cyclic structural matrices by admissibly permuting the rows to obtain a nonzero diagonal and then rescaling each row to make that diagonal equal to one.}
\end{lemma}

\begin{lemma}[lem:lacerda-distribution-entailment-equivalence]\label{lem:lacerda-distribution-entailment-equivalence}
\texttt{Theorem 1 of LacerdaSpirtesRamseyHoyer2008 states that the output of LiNG-\allowbreak{}D is an equivalence class under distribution entailment.}
\end{lemma}

\begin{lemma}[lem:kruskal-three-factor-uniqueness]\label{lem:kruskal-three-factor-uniqueness}
Let a three-way tensor have a decomposition into \(R\) nonzero rank-one terms with factor matrices of Kruskal ranks \(k_1,k_2,k_3\). If \(k_1+k_2+k_3\ge 2R+2\), then the tensor rank is \(R\) and the decomposition is unique up to a common permutation of the terms and reciprocal nonzero scalings within each term.
\end{lemma}

\begin{lemma}[lem:marcinkiewicz-cumulant-tail]\label{lem:marcinkiewicz-cumulant-tail}
\texttt{If a nondegenerate real random variable has moments of every order and only finitely many nonzero cumulants,\allowbreak{} then it is Gaussian. Consequently a nondegenerate non-\allowbreak{}Gaussian variable with moments of every order has infinitely many nonzero cumulants of order at least three.}
\end{lemma}
Apply the finite-cumulant (moment-sequence) form of Marcinkiewicz's theorem. If a probability law has moments of all orders and its moment-defined cumulants vanish above some finite order, the theorem concludes that the cumulant polynomial has degree at most two and that the law is Gaussian (degenerate cases included separately). This is the applicable formulation here: cumulants are finite partition polynomials in the assumed moments, so finite support of that sequence is exactly its premise. We do not infer that the characteristic function equals exp(P) from a formal Taylor series; exp(P) is only the familiar analytic-characteristic-function corollary of the cited finite-cumulant theorem. Thus a nondegenerate non-Gaussian law cannot have finitely many nonzero cumulants. Since its first two cumulants are finite, infinitely many of the nonzero cumulants must have order at least three.

\begin{lemma}[lem:carleman-moment-determinacy]\label{lem:carleman-moment-determinacy}
If a probability law on \(\mathbb R\) has finite moments \(m_k\) of all orders and \(\sum_{s=1}^{\infty}m_{2s}^{-1/(2s)}=\infty\), then the law is determined by its moment sequence.
\end{lemma}

\begin{lemma}[lem:veronese-lagrange-rank]\label{lem:veronese-lagrange-rank}
Let \(R\ge1\), let \(z_1,\ldots,z_R\) be nonzero pairwise nonproportional vectors in a finite-dimensional real vector space, and let \(s\ge R-1\). Then the columns \(z_1^{\otimes s},\ldots,z_R^{\otimes s}\) have Kruskal rank \(R\); equivalently, every subcollection is linearly independent.
\end{lemma}
\begin{proof}
It suffices to prove independence for an arbitrary subcollection \(z_1,\ldots,z_k\), where \(1\le k\le R\). For each ordered pair \(i\ne j\), nonproportionality supplies a linear functional \(\ell_{ij}\) such that \(\ell_{ij}(z_j)=0\) and \(\ell_{ij}(z_i)\ne0\): take a functional vanishing on the line spanned by \(z_j\) but not on \(z_i\). Choose also a linear functional \(g_i\) with \(g_i(z_i)\ne0\). The homogeneous polynomial
\[
 f_i(z)=g_i(z)^{s-k+1}\prod_{j\ne i}\ell_{ij}(z)
\]
has degree \(s\), is nonzero at \(z_i\), and vanishes at every \(z_j\) with \(j\ne i\). Every homogeneous polynomial of degree \(s\) is a linear functional on the symmetric tensor power, so there is \(L_i\) with \(L_i(z^{\otimes s})=f_i(z)\). If \(\sum_{j=1}^k a_jz_j^{\otimes s}=0\), applying \(L_i\) gives \(a_if_i(z_i)=0\), hence \(a_i=0\). This holds for every \(i\), proving independence of the arbitrary subcollection and therefore Kruskal rank \(R\).
\end{proof}

\begin{lemma}[lem:finite-moment-near-gaussian-witness]\label{lem:finite-moment-near-gaussian-witness}
For every integer \(K\ge3\) and every \(\rho>0\), there exists a centered, variance-one, non-Gaussian, moment-determinate probability law \(F\) on \(\mathbb R\) such that \(d_{\mathrm{TV}}(F,\mathcal N(0,1))<\rho\) and the first \(K\) raw moments, hence the cumulants through order \(K\), agree with those of \(\mathcal N(0,1)\).
\end{lemma}
\begin{proof}
Let \(\phi\) be the standard Gaussian density. Choose \(K+2\) pairwise disjoint bounded intervals \(I_0,\ldots,I_{K+1}\), each of positive Gaussian probability. The \((K+1)\)-by-\((K+2)\) matrix
\[
 A_{r\ell}=\int_{I_\ell}x^r\phi(x)\,dx,
 \qquad 0\le r\le K,
\]
has a nonzero null vector \(b=(b_0,\ldots,b_{K+1})\). Rescale it so that \(\max_\ell|b_\ell|=1\), and define the bounded, nonzero function \(h=\sum_\ell b_\ell\mathbf 1_{I_\ell}\). Then
\[
 \int x^rh(x)\phi(x)\,dx=0,
 \qquad 0\le r\le K. \tag{22}
\]
For \(0<t<1\), set \(f_t(x)=\phi(x)(1+th(x))\). The case \(r=0\) of (22) makes \(f_t\) integrate to one, and \(|h|\le1\) makes it strictly positive. Since \(h\) is nonzero on a set of positive Gaussian measure, \(f_t\ne\phi\). Equation (22) says that its first \(K\) moments, including mean zero and variance one, equal the Gaussian moments. Cumulants through order \(K\) are universal polynomials in the moments through the same order, so they also agree.

Moreover \(f_t\le2\phi\), whence its even moments satisfy \(m_{2s}(f_t)\le2(2s-1)!!\). It follows that \(m_{2s}(f_t)^{-1/(2s)}\) is bounded below by a positive constant times \(s^{-1/2}\), and therefore the Carleman series diverges. The cited Carleman criterion makes \(F_t\) moment determinate. Finally,
\[
 d_{\mathrm{TV}}(F_t,\mathcal N(0,1))
 =\frac t2\int|h(x)|\phi(x)\,dx\longrightarrow0
\]
as \(t\downarrow0\). Choosing \(t\) small enough proves the claim.
\end{proof}

\begin{theorem}[thm:heterogeneous-order-population-inverse-boundary]\label{thm:heterogeneous-order-population-inverse-boundary}
For a fixed normalized full-row-rank representation \(X=C\varepsilon\) with mutually independent, nondegenerate, non-Gaussian sources having finite absolute moments of every positive integer order and with pairwise nonproportional mixing columns, put
\[
 A_\nu=\{j:\operatorname{cum}_\nu(\varepsilon_j)\ne0\},
 \qquad R_\nu=|A_\nu|.
\]
The handle \(\mathcal H_{\mathrm{multi}}(P_X,n;K)\), run with exact population cumulants and the known source count, is a terminating pointwise population-oracle procedure: it processes \(\nu=2n-1,2n,\ldots\); for each processed order it solves the stated finite polynomial feasibility systems for \(R=0,\ldots,n\), selects the least feasible \(R\), and deduplicates the returned directions by exact projective equality until their union has cardinality \(n\). At every \(\nu\ge2n-1\), reshaping into blocks of degrees \((n-1,n-1,\nu-2n+2)\) makes the least feasible decomposition recover exactly \(\{[c_j]:j\in A_\nu\}\). The procedure stops after a finite representation-dependent order and returns
\[
 \mathcal D(P_X)=\{[c_1],\ldots,[c_n]\},
\]
without cross-order label alignment, and hence evaluates \(\mathcal V_x(C)\) and every \(\mathcal R_{xy}(C)\) from any representatives of those directions. For \(R_\nu\ge2\), recovery at a processed order has the ex-post adaptive certificate \(\nu\ge2R_\nu-1\) using block degrees \((R_\nu-1,R_\nu-1,\nu-2R_\nu+2)\); the universally executable schedule nevertheless begins at \(2n-1\). The cases \(R_\nu=0\) and \(R_\nu=1\) are respectively the zero-tensor and rank-one cases. This is neither a uniform finite-cutoff claim nor a finite-sample algorithm.

No finite cumulant cutoff is uniform over the qualitative all-finite-moments class. The negative conclusion holds even on the stricter subclass in which every source law is moment determinate. Already for \(p=n=2\), that stricter subclass contains two fixed separated direction and frontier targets and triangular sequences of observational laws whose \(N\)-sample total-variation distance tends to zero. Consequently neither the direction set nor the causal frontier admits a uniformly consistent estimator on that subclass, and no frontier-containing confidence-union sequence with uniform coverage strictly greater than one half can have uniformly vanishing Hausdorff error. The frontier map is also discontinuous at vanishing treatment loadings and deletion-rank boundaries. Thus a uniform root-\(N\) heterogeneous-order frontier law does not follow from the qualitative assumptions without quantitative non-Gaussianity, moment, Veronese-conditioning, angular, probe/eigen-gap, deletion-rank, and denominator margins. The existing common-order separated theorem retains its simultaneous coverage and root-\(N\) conclusion.
\end{theorem}
\begin{proof}
For each source and cumulant order define
\[
 \lambda_{j,\nu}=\operatorname{cum}_\nu(\varepsilon_j).
\]
The replaced assumption \(\mathrm{ass{:}moment\text{-}determinate\text{-}sources}\) states that \(\mathbb E|\varepsilon_j|^k<\infty\) for every \(j\in\mathsf S\) and every integer \(k\ge1\). Thus every raw moment and every cumulant used below exists. The joint-cumulant partition identity is
\[
 \operatorname{cum}(Z_1,\ldots,Z_\nu)
 =\sum_{\pi\in\mathfrak P_\nu}
 (|\pi|-1)!(-1)^{|\pi|-1}
 \prod_{B\in\pi}\mathbb E\!\left[\prod_{s\in B}Z_s\right]. \tag{20}
\]
This is a finite multilinear polynomial in joint moments. If the arguments split into two nonempty independent subfamilies, factorization of every moment in (20), followed by grouping partitions by the induced partitions on the two subfamilies, makes the alternating sum zero. Substituting \(X_i=\sum_j C_{ij}\varepsilon_j\), using multilinearity, and invoking \(\mathrm{ass{:}independent\text{-}sources}\) therefore gives the observable tensor identity
\[
 \mathcal T_\nu:=\operatorname{cum}_\nu(X)
 =\sum_{j=1}^n\lambda_{j,\nu}c_j^{\otimes\nu}
 =\sum_{j\in A_\nu}\lambda_{j,\nu}c_j^{\otimes\nu}. \tag{21}
\]

Fix \(\nu\ge2n-1\), put \(s=n-1\) and \(t=\nu-2n+2\), and reshape (21) into three blocks. Since \(n\ge p\ge2\), one has \(s\ge1\), and the order restriction gives \(t\ge1\). Thus
\[
 \widetilde{\mathcal T}_\nu
 =\sum_{j\in A_\nu}\lambda_{j,\nu}
 a_j\otimes b_j\otimes d_j,
 \qquad
 a_j=b_j=\operatorname{vec}(c_j^{\otimes s}),
 \qquad
 d_j=\operatorname{vec}(c_j^{\otimes t}). \tag{22}
\]
Let \(R=R_\nu\). Suppose first that \(R\ge2\). Because the columns of \(C\) are pairwise nonproportional and \(s=n-1\ge R-1\), \(\mathrm{lem{:}veronese\text{-}lagrange\text{-}rank}\) gives Kruskal rank \(R\) to each of the first two factor matrices in (22). Every pair of third-mode columns is linearly independent. Indeed, if \(z^{\otimes t}\) and \(w^{\otimes t}\) were proportional, then contraction against \(f^{\otimes t}\), for any linear functional \(f\) vanishing on \(w\), would give \(f(z)^t=0\). Every functional vanishing on \(w\) would then vanish on \(z\), implying \(z\in\operatorname{span}(w)\), contrary to distinct directions. Hence the third factor matrix has Kruskal rank at least two, and
\[
 k_1+k_2+k_3\ge R+R+2=2R+2. \tag{23}
\]
By \(\mathrm{lem{:}kruskal\text{-}three\text{-}factor\text{-}uniqueness}\), the decomposition (22) is unique up to a common permutation and reciprocal nonzero scalings.

The length \(R\) is minimal. Write \(A=[a_j]_{j\in A_\nu}\), \(B=[b_j]_{j\in A_\nu}\), and \(K=[b_j\otimes d_j]_{j\in A_\nu}\). The first two matrices have full column rank. If \(\sum_j q_jb_j\otimes d_j=0\), apply the row of a left inverse of \(B\) selecting column \(j\) to the first tensor component; this yields \(q_jd_j=0\), and \(d_j\ne0\) yields \(q_j=0\). Thus \(K\) also has full column rank. The first matricization of (22) is
\[
 A\operatorname{diag}(\lambda_{j,\nu})K^\top, \tag{24}
\]
whose matrix rank is \(R\), because both outer factors have column rank \(R\) and every \(\lambda_{j,\nu}\), \(j\in A_\nu\), is nonzero by definition. A sum of fewer than \(R\) rank-one tensors has first-matricization rank below \(R\), contradicting (24). Thus the minimum number of terms is exactly \(R\), and Kruskal uniqueness implies that every least feasible decomposition in \(\mathrm{def{:}multi\text{-}order\text{-}handle}\) returns precisely \(\{[c_j]:j\in A_\nu\}\).

This minimization is exposed as a finite population-oracle operation rather than an opaque factorization. At fixed \(\nu\) and \(R\), its variables are \(\gamma_1,\ldots,\gamma_R\in\mathbb R\) and \(z_1,\ldots,z_R\in\mathbb R^p\). Its constraints are the finitely many coordinate equations
\[
 \widetilde{\mathcal T}_\nu
 =\sum_{h=1}^{R}\gamma_h
 \operatorname{vec}(z_h^{\otimes(n-1)})\otimes
 \operatorname{vec}(z_h^{\otimes(n-1)})\otimes
 \operatorname{vec}(z_h^{\otimes(\nu-2n+2)}), \tag{25}
\]
together with the polynomial inequalities \(\gamma_h^2>0\) and \(\lVert z_h\rVert_2^2>0\). The handle asks its exact feasibility oracle about (25) for the finite list \(R=0,\ldots,n\) and chooses the least feasible value. Exact projective equality of two returned nonzero vectors is tested by the vanishing of all two-by-two minors of the matrix having those vectors as columns. Thus every operation at a processed order is explicit. The population law supplies exact cumulants; no claim is made that these oracle operations constitute a finite-sample numerical algorithm.

If \(R=0\), (21) is the zero tensor and the empty decomposition is the least feasible output. If \(R=1\), it is a nonzero rank-one tensor. The projective Veronese map is injective: the contraction argument used after (22) shows that proportional nonzero tensor powers imply proportional base vectors. Hence the unique active direction is recovered. This covers all active ranks.

The same proof gives the adaptive ex-post certificate. If \(R_\nu\ge2\) and \(\nu\ge2R_\nu-1\), use degrees
\[
 (R_\nu-1,R_\nu-1,\nu-2R_\nu+2). \tag{26}
\]
The first two degrees are at least \(R_\nu-1\), and the third is positive. The first two Kruskal ranks are therefore \(R_\nu\), the third is at least two, and (23)--(24) prove uniqueness and minimality. Because \(R_\nu\) is learned only after solving the order-\(\nu\) systems, (26) is an ex-post certificate. The universally executable schedule in \(\mathrm{def{:}multi\text{-}order\text{-}handle}\) uses the known \(n\) and begins at \(2n-1\); it does not presuppose \(R_\nu\).

By \(\mathrm{lem{:}marcinkiewicz\text{-}cumulant\text{-}tail}\), the all-finite-moments, nondegeneracy, and non-Gaussianity assumptions imply that every source has infinitely many nonzero cumulants of order at least three. This set of orders is unbounded, so for each \(j\) there is a finite \(\nu_j\ge2n-1\) with \(\lambda_{j,\nu_j}\ne0\). Consequently
\[
 K_{C,\varepsilon}=\max_{1\le j\le n}\nu_j<\infty. \tag{27}
\]
By the orderwise recovery proved above, processing \(2n-1,2n,\ldots,K_{C,\varepsilon}\) and unioning the outputs includes every \([c_j]\). Pairwise nonproportionality prevents erroneous projective deduplication, and no direction outside this set is returned. The handle therefore stops no later than (27), when the union first has cardinality \(n\), and returns
\[
 \{[c_1],\ldots,[c_n]\}=\mathcal D(P_X), \tag{28}
\]
where the equality uses \(\mathrm{thm{:}sharp\text{-}effect\text{-}frontier}\). No cross-order labeling is needed. Choosing an arbitrary nonzero representative of each recovered direction produces \(C'=C\Pi D\); \(\mathrm{lem{:}frontier\text{-}monomial\text{-}invariance}\) then yields the same \(\mathcal V_x\) and \(\mathcal R_{xy}\). The sharp-effect theorem, rather than a definition, proves that these evaluated sets equal the causal law-fiber response targets. This completes the positive result using only all finite moments, nondegeneracy, non-Gaussianity, independence, normalized full-row-rank pairwise-nonproportional mixing, and known \(n\); moment determinacy has not been used.

We next prove that no cumulant cutoff is uniform even on the stricter moment-determinate subclass. Fix any finite \(K\ge3\) and \(\rho>0\). By \(\mathrm{lem{:}finite\text{-}moment\text{-}near\text{-}gaussian\text{-}witness}\), there is a centered, variance-one, non-Gaussian, moment-determinate law \(F\) within total variation \(\rho\) of \(G=\mathcal N(0,1)\), whose moments through order \(K\) agree with those of \(G\). In particular, this witness has moments of every order and belongs to the positive theorem's all-finite-moments class. Give two sources independent law \(F\), fix \(0<\gamma<\pi/4\), and set
\[
 C_0=I_2,
 \qquad
 C_1=
 \begin{pmatrix}
 \cos\gamma&-\sin\gamma\\
 \sin\gamma&\cos\gamma
 \end{pmatrix}. \tag{29}
\]
Both matrices are normalized and full rank and have nonzero pairwise nonproportional columns. Independence makes all source joint moments of total degree at most \(K\) equal the corresponding moments under \(G^{\otimes2}\). Since \(C_1\) is orthogonal and \(G^{\otimes2}\) is orthogonally invariant, the two observed laws in (29) have identical moments, hence identical cumulants, through order \(K\), while their projective direction sets differ. Because \(K\) was arbitrary, matching through any proposed cutoff delays all informative cumulants beyond that cutoff. The Marcinkiewicz lemma nevertheless guarantees a nonzero cumulant at some later order for each witness source. Therefore no fixed finite cumulant cutoff recovers the directions throughout even the moment-determinate subclass.

For the statistical boundary, for each \(N\) choose the preceding moment-determinate witness with
\[
 d_{\mathrm{TV}}(F_N,G)\le\frac{1}{4N^2}, \tag{30}
\]
and let \(P_{k,N}\) be the law generated by \(C_k\) and two independent \(F_N\) sources. Total variation contracts under measurable pushforward because the inverse image of each measurable event is measurable. Also, a telescoping expansion of product signed measures and the triangle inequality give
\[
 d_{\mathrm{TV}}(\mu^{\otimes m},\nu^{\otimes m})
 \le m\,d_{\mathrm{TV}}(\mu,\nu). \tag{31}
\]
Each observed law is within \(2d_{\mathrm{TV}}(F_N,G)\) of the common law \(G^{\otimes2}\), the latter remaining unchanged by either orthogonal matrix in (29). The triangle inequality, (30), and then (31) imply
\[
 d_{\mathrm{TV}}(P_{0,N},P_{1,N})\le N^{-2},
 \qquad
 d_{\mathrm{TV}}(P_{0,N}^{\otimes N},P_{1,N}^{\otimes N})\le N^{-1}\longrightarrow0. \tag{32}
\]

The two projective direction sets in (29) have a fixed positive Hausdorff separation in any fixed metric compatible with real projective space. A uniformly consistent direction-set estimator would yield a test by selecting the closer of these two fixed targets, with both errors tending to zero. For every test event \(E\), however, the definition of total variation gives
\[
 P_{0,N}^{\otimes N}(E^c)+P_{1,N}^{\otimes N}(E)
 \ge1-d_{\mathrm{TV}}(P_{0,N}^{\otimes N},P_{1,N}^{\otimes N}), \tag{33}
\]
which tends to one by (32), a contradiction.

The corresponding causal frontiers are fixed and separated. Directly applying \(\mathrm{def{:}admissible\text{-}source\text{-}set}\) at \(x=1\), \(y=2\), gives
\[
 \mathcal R_{12}(C_0)=\{0\},
 \qquad
 \mathcal R_{12}(C_1)=\{\tan\gamma,-\cot\gamma\}. \tag{34}
\]
Thus the nearest-target argument also rules out a uniformly consistent frontier estimator. More strongly, let \(\mathcal U_N\) be a random compact set with uniform frontier coverage at least \(1-\beta\), where \(\beta<1/2\). Under \(P_{0,N}^{\otimes N}\), coverage of the first target has probability at least \(1-\beta\). Coverage of the second target has probability at least \(1-\beta\) under \(P_{1,N}^{\otimes N}\), and (32) transfers that event to \(P_{0,N}^{\otimes N}\) with probability at least \(1-\beta-o(1)\). Their intersection consequently has probability at least \(1-2\beta-o(1)>0\). On the intersection, since \(0\in\mathcal R_{12}(C_0)\) and \(-\cot\gamma\in\mathcal U_N\),
\[
 d_H(\mathcal U_N,\mathcal R_{12}(C_0))\ge\cot\gamma. \tag{35}
\]
Hence no uniformly honest frontier-containing confidence union with coverage strictly above one half has uniformly vanishing Hausdorff error, even on this moment-determinate subclass.

Finally, two explicit families establish the deterministic discontinuity boundaries. Column normalization may be applied because \(\mathrm{lem{:}frontier\text{-}monomial\text{-}invariance}\) preserves the frontiers. For
\[
 D_t=\begin{pmatrix}t&1\\1&1\end{pmatrix},
\]
the matrices are full rank and have nonzero nonproportional columns for all sufficiently small \(|t|\) with \(t\ne1\), and deletion-rank calculation gives
\[
 \mathcal R_{12}(D_t)=\{1/t,1\}\quad(t\ne0),
 \qquad
 \mathcal R_{12}(D_0)=\{1\}. \tag{36}
\]
This proves discontinuity at a vanishing treatment loading. For
\[
 E_t=
 \begin{pmatrix}
 1&1&0&1\\
 1&t&0&0\\
 0&0&1&1
 \end{pmatrix},
\]
the matrices are full row rank with nonzero nonproportional columns for all sufficiently small \(|t|\). Deleting row one and the indicated column shows
\[
 \mathcal R_{12}(E_0)=\{0\},
 \qquad
 \mathcal R_{12}(E_t)=\{0,t,1\}\quad(t\ne0). \tag{37}
\]
Indeed, at \(t=0\) only columns two and four are admissible and both ratios are zero; at \(t\ne0\), columns one, two, and four are admissible and give ratios \(1,t,0\). This proves discontinuity at a deletion-rank boundary.

Equations (32), (36), and (37) show that the qualitative all-finite-moments assumptions alone supply neither a uniform stopping order nor a uniform statistical modulus. Quantitative non-Gaussianity, moment, Veronese-conditioning, angular, probe/eigen-gap, denominator, and deletion-rank margins therefore cannot be omitted from a uniform root-\(N\) result. On the declared separated common-order iid class, \(\mathrm{thm{:}uniform\text{-}confidence\text{-}frontier}\) supplies simultaneous response-vector and scalar coverage with \(O(N^{-1/2})\) Hausdorff error. This completes the pointwise population-oracle construction and the nonuniform inference boundary without strengthening any assumption.
\end{proof}
\par\smallskip\noindent \textit{Justification.} The theorem gives a terminating pointwise population-oracle schedule under all finite source moments: it starts at \(2n-1\), solves explicit finite polynomial systems, unions source-specific informative orders, and stops at a representation-dependent order. Moment determinacy is not used positively; the negative nonuniformity result nevertheless holds even for moment-determinate witnesses. Its negative target is precise: it rules out a putative extension that fixes a maximum cumulant order in advance and uniformly recovers the direction/frontier, or supplies a uniformly shrinking frontier confidence union, over the unseparated qualitative all-finite-moments class. It does not refute the model-specific PodosinnikovaEtAl2019 SDP guarantee. The accepted-bank bivariate arrow theorems identify one axis-conditioned arrow at lower fixed generic truncations, whereas this theorem recovers all directions in arbitrary observed dimension and maps them to a cyclic response frontier, so the results are incomparable. \textit{Gap.} PodosinnikovaEtAl2019, Algorithm 1, estimates a cumulant-generating-function-Hessian subspace and applies an SDP relaxation; its abstract gives high-probability recovery of a mixing component when the mixing components are uniformly random on the sphere and \(k<(2-\epsilon)p\log p\). The negative half here neither weakens nor contradicts that model-specific result. Rather, it supplies the missing boundary for the different theorem-level target of a heterogeneous-order extension with a predeclared finite cutoff and uniform direction/frontier consistency or uniformly shrinking frontier confidence sets across all finite-moment non-Gaussian sources without quantitative margins. WangSeigal2026 supplies full-law identifiability and a coupled second/fourth-cumulant algorithm, while the accepted-bank bivariate arrow theorems reach fixed generic truncation \(2n-2\), and at most \(2n-3\) only for \(m\ge3\), equivalently \(n\ge5\). None supplies this representation-dependent union of source-specific informative orders under an all-finite-moments premise, its deletion-rank frontier evaluation, or the impossibility of uniform stopping and shrinking inference even on the stricter moment-determinate subclass. \textit{Consumer.} Analysts with exact population cumulants, known source count, and all finite source moments can evaluate all directions and the whole frontier without a common informative order. The negative half says exactly which extrapolation is unavailable: not the cited random-spherical SDP guarantee, but any unseparated heterogeneous-order procedure promised a fixed finite cutoff together with uniform recovery or a uniformly shrinking frontier confidence union.

\begin{lemma}[lem:dai-irreducible-oica-uniqueness]\label{lem:dai-irreducible-oica-uniqueness}
Let \(C\varepsilon\) and \(A^{\widetilde{\mathcal M}}\xi^{\widetilde{\mathcal M}}\) be two irreducible independent-component representations of the same \(p\)-dimensional law, where irreducible means that the mixing matrix has neither zero columns nor proportional column pairs. If \(C\in\mathbb R^{p\times n}\), \(A^{\widetilde{\mathcal M}}\in\mathbb R^{p\times m_{\widetilde{\mathcal M}}}\), the source coordinates are mutually independent and nondegenerate, and every coordinate of \(\varepsilon\) is non-Gaussian, then \(m_{\widetilde{\mathcal M}}=n\), every column of one mixing matrix is proportional to a unique column of the other, and every coordinate of \(\xi^{\widetilde{\mathcal M}}\) is non-Gaussian. Equivalently, there are a permutation matrix \(\Pi\) and a nonsingular diagonal matrix \(D\) such that
\[
 A^{\widetilde{\mathcal M}}=C\Pi D.
\]
\end{lemma}

\begin{lemma}[lem:law-hard-intervention-response]\label{lem:law-hard-intervention-response}
For every \(\widetilde{\mathcal M}\in\mathfrak F_x^{\mathrm{law}}(P_X)\), the response vector is well defined and
\[
 \rho_x(\widetilde{\mathcal M})
 =\left(
 \frac{((Q^{\widetilde{\mathcal M}})^{-1})_{ix}}
 {((Q^{\widetilde{\mathcal M}})^{-1})_{xx}}
 \right)_{i\in\mathsf O\setminus\{x\}}.
\]
\end{lemma}
\begin{proof}
Write \(Q=Q^{\widetilde{\mathcal M}}\), \(H=H^{\widetilde{\mathcal M}}\), and \(\xi=\xi^{\widetilde{\mathcal M}}\). Membership in the law fiber gives \(\det(Q)\ne0\) and \(\det(Q_{-x,-x})\ne0\). The diagonal cofactor formula therefore gives
\[
 (Q^{-1})_{xx}=\frac{\det(Q_{-x,-x})}{\det(Q)}\ne0. \tag{1}
\]
After the perfect intervention replaces the \(x\)-th equation by \(V_x=t\), the remaining equations are
\[
 Q_{-x,-x}V_{-x}=H_{-x,:}\xi-Q_{-x,x}t. \tag{2}
\]
The postintervention determinant condition makes (2) uniquely solvable, and differentiation with the disturbances held fixed yields
\[
 \frac{\partial V_{-x}}{\partial t}=-Q_{-x,-x}^{-1}Q_{-x,x}. \tag{3}
\]
On the other hand, the rows different from \(x\) in \(QQ^{-1}e_x=e_x\) give
\[
 Q_{-x,-x}(Q^{-1})_{-x,x}+Q_{-x,x}(Q^{-1})_{xx}=0. \tag{4}
\]
Multiplying (4) by \(Q_{-x,-x}^{-1}\) and dividing by the nonzero quantity in (1) shows that the right-hand side of (3) equals \((Q^{-1})_{-x,x}/(Q^{-1})_{xx}\). Restricting to the observed coordinates different from \(x\) and invoking the definition of \(\rho_x\) proves the claim.
\end{proof}

\begin{lemma}[lem:frontier-monomial-invariance]\label{lem:frontier-monomial-invariance}
Let \(C'=C\Pi D\), where \(\Pi e_k=e_{\pi(k)}\) is a permutation matrix and \(D=\operatorname{diag}(d_1,\ldots,d_n)\) has \(d_k\ne0\) for every \(k\). Then
\[
 k\in J_x(C')\quad\Longleftrightarrow\quad \pi(k)\in J_x(C),
\]
and consequently \(\mathcal V_x(C')=\mathcal V_x(C)\) and \(\mathcal R_{xy}(C')=\mathcal R_{xy}(C)\).
\end{lemma}
\begin{proof}
The \(k\)-th column of \(C'\) is \(d_kc_{\pi(k)}\), so
\[
 C'_{xk}=d_kC_{x,\pi(k)}. \tag{5}
\]
After deleting column \(k\), the remaining columns are a permutation and nonzero diagonal rescaling of the columns of \(C\) other than \(\pi(k)\). Hence there are an \((n-1)\)-dimensional permutation matrix \(\Pi^{(k)}\) and a nonsingular diagonal matrix \(D^{(k)}\) such that
\[
 C'_{-x,-k}=C_{-x,-\pi(k)}\Pi^{(k)}D^{(k)}. \tag{6}
\]
Right multiplication by the two invertible matrices in (6) preserves rank. Equations (5)--(6) and the definition of \(J_x\) prove the displayed equivalence. Moreover,
\[
 \frac{C'_{ik}}{C'_{xk}}
 =\frac{d_kC_{i,\pi(k)}}{d_kC_{x,\pi(k)}}
 =\frac{C_{i,\pi(k)}}{C_{x,\pi(k)}} \tag{7}
\]
for every admissible \(k\) and every observed \(i\). Since \(\pi\) is bijective, (7) proves both frontier equalities from their definitions.
\end{proof}

\begin{lemma}[lem:monomial-law-canonicalization]\label{lem:monomial-law-canonicalization}
Let \(\widetilde{\mathcal M}\in\mathfrak F_x^{\mathrm{law}}(P_X)\) and suppose \(m_{\widetilde{\mathcal M}}=n\) and \(A^{\widetilde{\mathcal M}}=C\Pi D\), where \(\Pi e_k=e_{\pi(k)}\) and \(D\) is nonsingular diagonal. Define
\[
 H^C=H^{\widetilde{\mathcal M}}D^{-1}\Pi^\top,
 \qquad
 \mathcal M^C=(Q^{\widetilde{\mathcal M}},H^C).
\]
Then \(\mathcal M^C\in\mathfrak F_x(C,\varepsilon)\). For every \(i\in\mathsf O\), if \(s_{\widetilde{\mathcal M}}^{\mathrm{law}}(i)=k\), then \(s_{\mathcal M^C}(i)=\pi(k)=\bar s_{\widetilde{\mathcal M}}(i)\). The inverse structural matrix, its intervention minor, and every equilibrium response ratio are unchanged by this canonicalization.
\end{lemma}
\begin{proof}
The matrix \(D^{-1}\Pi^\top\) is monomial. Because the product of two monomial matrices is monomial, the source-assignment condition for \(H^{\widetilde{\mathcal M}}\) implies that \(H^C\) is monomial. The matrix \(Q^{\widetilde{\mathcal M}}\) is not changed, so its unit diagonal, its nonzero determinant, and the nonzero determinant of its \((-x,-x)\) minor are inherited from membership in the law fiber. The observed mixing rows satisfy
\[
 \begin{aligned}
 \bigl[(Q^{\widetilde{\mathcal M}})^{-1}H^C\bigr]_{\mathsf O,:}
 &=A^{\widetilde{\mathcal M}}D^{-1}\Pi^\top\\
 &=C\Pi DD^{-1}\Pi^\top=C. \tag{8}
 \end{aligned}
\]
Thus all member conditions in the fixed-law completion-fiber definition hold, and \(\mathcal M^C\in\mathfrak F_x(C,\varepsilon)\).

If row \(i\) of \(H^{\widetilde{\mathcal M}}\) has its unique nonzero in column \(k\), then right multiplication by \(D^{-1}\Pi^\top\) moves that nonzero to column \(\pi(k)\) and rescales it by a nonzero number. Hence \(s_{\mathcal M^C}(i)=\pi(k)\). The alignment equation says that column \(k\) of \(A^{\widetilde{\mathcal M}}\) is \(d_kc_{\pi(k)}\); distinct reference directions make \(\pi(k)\) unique, so it is also \(\bar s_{\widetilde{\mathcal M}}(i)\). Finally, canonicalization leaves \(Q^{\widetilde{\mathcal M}}\) unchanged. The hard-intervention response formula therefore shows that its inverse-matrix ratios, its intervention minor, and all response coordinates are unchanged.
\end{proof}

\begin{lemma}[lem:law-fiber-response-rank]\label{lem:law-fiber-response-rank}
For every \(\widetilde{\mathcal M}\in\mathfrak F_x^{\mathrm{law}}(P_X)\), one has \(m_{\widetilde{\mathcal M}}=n\), the aligned source \(j=\bar s_{\widetilde{\mathcal M}}(x)\) belongs to \(J_x(C)\), and
\[
 \rho_x(\widetilde{\mathcal M})
 =\left(\frac{C_{ij}}{C_{xj}}\right)_{i\in\mathsf O\setminus\{x\}}
 \in\mathcal V_x(C).
\]
In particular, \(C_{xj}\ne0\) and \(\operatorname{rank}(C_{-x,-j})=p-1\).
\end{lemma}
\begin{proof}
The reference representation \((C,\varepsilon)\) is irreducible because its columns are nonzero and pairwise nonproportional. Its sources are mutually independent, nondegenerate, and non-Gaussian by the corresponding population assumptions. The representation \((A^{\widetilde{\mathcal M}},\xi^{\widetilde{\mathcal M}})\) is irreducible by the law-fiber definition, and it has the same observed law. The cited irreducible-OICA uniqueness lemma therefore gives
\[
 m_{\widetilde{\mathcal M}}=n,
 \qquad
 A^{\widetilde{\mathcal M}}=C\Pi D \tag{9}
\]
for a permutation matrix \(\Pi\) and a nonsingular diagonal matrix \(D\).

Apply the monomial-canonicalization lemma to (9). It returns \(\mathcal M^C\in\mathfrak F_x(C,\varepsilon)\) whose source assigned to equation \(x\) is
\[
 s_{\mathcal M^C}(x)=\bar s_{\widetilde{\mathcal M}}(x)=j. \tag{10}
\]
The established necessary deletion-rank lemma applied to \(\mathcal M^C\) yields \(j\in J_x(C)\); expanding the definition of \(J_x(C)\) gives the final two asserted conditions.

It remains to identify the whole vector. The established intervention-ratio lemma is stated for an arbitrary observed outcome distinct from \(x\), so instantiate it at each \(i\in\mathsf O\setminus\{x\}\). For the fixed-fiber completion \(\mathcal M^C\) it gives
\[
 \frac{((Q^{\widetilde{\mathcal M}})^{-1})_{ix}}
 {((Q^{\widetilde{\mathcal M}})^{-1})_{xx}}
 =\frac{C_{ij}}{C_{xj}}, \tag{11}
\]
because canonicalization leaves \(Q^{\widetilde{\mathcal M}}\) unchanged. Combining (11) with the hard-intervention response formula proves the displayed vector equality. Its membership in \(\mathcal V_x(C)\) is then exactly the response-frontier definition.
\end{proof}
\par\smallskip\noindent \textit{Justification.} This lemma is the law-fiber necessity bridge from published ICA uniqueness to causal response and deletion rank. \textit{Gap.} Neither ICA uniqueness nor fixed-representation intervention algebra alone proves necessity for arbitrary same-law completions. \textit{Consumer.} It certifies that an arbitrary legal completion cannot escape the observable vector frontier.

\begin{lemma}[lem:numerical-common-order-local-inverse]\label{lem:numerical-common-order-local-inverse}
Let \(C,C'\in\mathbb R^{p\times n}\) and \(\lambda,\lambda'\in\mathbb R^n\). Suppose \(d\ge1\), \(q\ge1\), \(r=2d+q\), and \(n\le\binom{p+d-1}{d}\). For every \(j\in\mathsf S\), suppose
\[
 \lVert c_j\rVert_2=\lVert c'_j\rVert_2=1,
 \qquad
 \kappa\le |\lambda_j|,|\lambda'_j|\le\Lambda,
 \qquad
 u^\top c_j,u^\top c'_j\ge\sigma.
\]
Suppose also that both lifted matrices \([c_1^{\otimes d}\;\cdots\;c_n^{\otimes d}]\) and \([(c'_1)^{\otimes d}\;\cdots\;(c'_n)^{\otimes d}]\) have least singular value at least \(\sigma\), and that both ratio collections
\[
 \left\{\frac{v^\top c_j}{u^\top c_j}:j\in\mathsf S\right\},
 \qquad
 \left\{\frac{v^\top c'_j}{u^\top c'_j}:j\in\mathsf S\right\}
\]
have pairwise gaps at least \(\sigma\). Define
\[
 \mathcal T(C,\lambda)=\sum_{j=1}^n\lambda_jc_j^{\otimes r},
 \qquad
 \mathcal T(C',\lambda')=\sum_{j=1}^n\lambda'_j(c'_j)^{\otimes r}.
\]
If \(\lVert\mathcal T(C',\lambda')-\mathcal T(C,\lambda)\rVert_F<\epsilon_0\), then there is a permutation \(\pi\) of \(\mathsf S\) such that
\[
 \lVert C'_{\pi}-C\rVert_F
 \le L_C\lVert\mathcal T(C',\lambda')-\mathcal T(C,\lambda)\rVert_F.
\]
In particular, this bound applies to two numerical factor pairs in the feasible set of the separated tensor minimum-distance fit; no source law for the fitted pair is required.
\end{lemma}
\begin{proof}
Put
\[
 e=\lVert\mathcal T(C',\lambda')-\mathcal T(C,\lambda)\rVert_F.
\]
For a unit vector \(w\in\mathbb R^p\), contract the last \(q\) modes against \(u^{\otimes(q-1)}\otimes w\), and matricize the other \(2d\) modes in two blocks of degree \(d\). By multilinearity, the resulting matrices are
\[
 M_w=V_dD_wV_d^\top,
 \qquad
 D_w=\operatorname{diag}\!\left(
 \lambda_j(u^\top c_j)^{q-1}(w^\top c_j)
 \right)_{j=1}^n,
 \tag{1}
\]
and analogously \(M'_w=V'_dD'_w(V'_d)^\top\). Contraction against unit vectors has operator norm one for the Frobenius norm, matricization preserves that norm, and orthogonal compression cannot increase it. Hence
\[
 \lVert M'_w-M_w\rVert_{\mathrm{op}}
 \le \lVert M'_w-M_w\rVert_F\le e.
 \tag{2}
\]

Let \(U\in\mathbb R^{p^d\times n}\) have orthonormal columns spanning \(\operatorname{range}(V_d)\), and put
\[
 S=U^\top V_d,
 \qquad
 A_w=U^\top M_wU=SD_wS^\top.
 \tag{3}
\]
The lifted-rank condition gives \(s_{\min}(S)=s_{\min}(V_d)\ge\sigma\). Unit normalization gives \(\lVert S\rVert_{\mathrm{op}}\le\lVert V_d\rVert_F=\sqrt n\). When \(w=u\), every diagonal entry of \(D_u\) has magnitude at least \(\kappa\sigma^q\), so the product inequality for least singular values yields
\[
 s_{\min}(A_u)
 \ge s_{\min}(S)^2s_{\min}(D_u)
 \ge\kappa\sigma^{q+2}=\eta,
 \qquad
 \lVert A_u^{-1}\rVert_{\mathrm{op}}\le\eta^{-1}.
 \tag{4}
\]
For every unit \(w\), the upper cumulant bound and unit normalization similarly give
\[
 \lVert A_w\rVert_{\mathrm{op}}
 \le \lVert S\rVert_{\mathrm{op}}^2\lVert D_w\rVert_{\mathrm{op}}
 \le n\Lambda.
 \tag{5}
\]

Compress the primed contractions with the same \(U\): let \(A'_w=U^\top M'_wU\). Equation (2) implies \(\lVert A'_w-A_w\rVert_{\mathrm{op}}\le e\). Since \(e<\epsilon_0\le\eta/2\), (4) and the Neumann-series identity give
\[
 \lVert(A'_u)^{-1}\rVert_{\mathrm{op}}
 \le\frac{\lVert A_u^{-1}\rVert_{\mathrm{op}}}
 {1-\lVert A_u^{-1}\rVert_{\mathrm{op}}\lVert A'_u-A_u\rVert_{\mathrm{op}}}
 \le2\eta^{-1}.
 \tag{6}
\]
Writing \(S'=U^\top V'_d\), one has \(A'_u=S'D'_u(S')^\top\). The diagonal matrix \(D'_u\) is nonsingular by the primed cumulant and probe margins. Since \(A'_u\) is nonsingular by (6), the square matrix \(S'\) is nonsingular. Consequently, for
\[
 G_w=A_wA_u^{-1},
 \qquad
 G'_w=A'_w(A'_u)^{-1},
\]
equations (1)--(3) give the simultaneous diagonalizations
\[
 G_w=S\operatorname{diag}\!\left(\frac{w^\top c_j}{u^\top c_j}\right)S^{-1},
 \qquad
 G'_w=S'\operatorname{diag}\!\left(\frac{w^\top c'_j}{u^\top c'_j}\right)(S')^{-1}.
 \tag{7}
\]
The inverse identity
\[
 (A'_u)^{-1}-A_u^{-1}
 =(A'_u)^{-1}(A_u-A'_u)A_u^{-1},
\]
together with (2), (4)--(6), gives, for every unit \(w\),
\[
 \lVert G'_w-G_w\rVert_{\mathrm{op}}
 \le\left(\frac2\eta+\frac{2n\Lambda}{\eta^2}\right)e
 =he.
 \tag{8}
\]

For \(w=v\), the eigenvalues of \(G_v\) are real and separated by at least \(\sigma\). Around each eigenvalue draw a circle in \(\mathbb C\) of radius \(g=\sigma/3\). Equation (7) implies that the resolvent of the complexification of \(G_v\) on every such circle is bounded by
\[
 \lVert(zI-G_v)^{-1}\rVert_{\mathrm{op}}
 \le\frac{\lVert S\rVert_{\mathrm{op}}\lVert S^{-1}\rVert_{\mathrm{op}}}{g}
 \le\frac\chi g.
 \tag{9}
\]
Because \(e<\epsilon_0\le\sigma/(6\chi h)\), equations (8)--(9) imply
\[
 \lVert(zI-G_v)^{-1}(G'_v-G_v)\rVert_{\mathrm{op}}<\frac12.
\]
Thus the Neumann series keeps every circle in the resolvent set along the path \(G_v+t(G'_v-G_v)\), \(0\le t\le1\). The Riesz projector rank is integer valued and continuous along this path, so every circle contains exactly one eigenvalue of \(G'_v\), counted algebraically. Label that eigenvalue by \(\pi(j)\). The resolvent identity, followed by integration around the circle, yields
\[
 \lVert P'_{\pi(j)}-P_j\rVert_{\mathrm{op}}
 \le g\left(\frac{2\chi}{g}\right)he\left(\frac\chi g\right)
 =\frac{6\chi^2he}{\sigma}.
 \tag{10}
\]
Moreover \(\lVert P_j\rVert_{\mathrm{op}}\le\chi\), and (10) together with \(e<\sigma/(6\chi h)\) gives \(\lVert P'_{\pi(j)}\rVert_{\mathrm{op}}\le2\chi\). Since the \(n\) disjoint circles each carry one eigenvalue, \(\pi\) is a permutation.

For the \(i\)-th coordinate vector \(e_i\in\mathbb R^p\), define
\[
 z_{ij}=\operatorname{tr}(G_{e_i}P_j),
 \qquad
 z'_{i,\pi(j)}=\operatorname{tr}(G'_{e_i}P'_{\pi(j)}).
\]
The simultaneous diagonalizations in (7), and simplicity of the \(v\)-pencil spectrum, imply
\[
 z_j=\frac{c_j}{u^\top c_j},
 \qquad
 z'_{\pi(j)}=\frac{c'_{\pi(j)}}{u^\top c'_{\pi(j)}}.
 \tag{11}
\]
Using \(|\operatorname{tr}(R)|\le n\lVert R\rVert_{\mathrm{op}}\), (5), (8), and (10), one obtains for every \(i,j\)
\[
 \begin{aligned}
 |z'_{i,\pi(j)}-z_{ij}|
 &\le n\left(
 \lVert G'_{e_i}-G_{e_i}\rVert_{\mathrm{op}}
 \lVert P'_{\pi(j)}\rVert_{\mathrm{op}}
 +\lVert G_{e_i}\rVert_{\mathrm{op}}
 \lVert P'_{\pi(j)}-P_j\rVert_{\mathrm{op}}
 \right)\\
 &\le nh\left(2\chi+\frac{6n\Lambda\chi^2}{\eta\sigma}\right)e
 =L_ze.
 \end{aligned}
 \tag{12}
\]
Hence \(\lVert z'_{\pi(j)}-z_j\rVert_2\le\sqrt pL_ze\). The probe margins orient both vectors positively, and unit normalization in (11) gives \(\lVert z_j\rVert_2,\lVert z'_{\pi(j)}\rVert_2\ge1\). For nonzero vectors \(z,z'\), the triangle inequality gives
\[
 \left\lVert\frac z{\lVert z\rVert_2}-
 \frac{z'}{\lVert z'\rVert_2}\right\rVert_2
 \le\frac{2\lVert z-z'\rVert_2}
 {\min\{\lVert z\rVert_2,\lVert z'\rVert_2\}}.
\]
Applying this inequality to (11)--(12) gives
\[
 \lVert c'_{\pi(j)}-c_j\rVert_2\le2\sqrt pL_ze.
\]
Squaring, summing over \(j\), and using \(L_C=2\sqrt{np}L_z\) proves
\[
 \lVert C'_{\pi}-C\rVert_F\le L_Ce.
\]
Every step used only the numerical hypotheses displayed in the statement, so the conclusion applies to the numerical feasible set of the minimum-distance estimator.
\end{proof}

\begin{lemma}[lem:shared-kernel-matching-certificate]\label{lem:shared-kernel-matching-certificate}
Suppose \(\operatorname{rank}(C)=p\), and run the shared-kernel and shared-matching construction in \(\operatorname{Cert}(C,x)\). Then \(K_x=(W_0)_{:,\{x\}\cup\mathsf L}\) is a basis matrix for \(\ker(C_{-x,:})\), and
\[
 j\in J_x(C)
 \quad\Longleftrightarrow\quad
 C_{xj}\ne0\ \text{ and the \(j\)-th row of \(K_x\) is nonzero}.
\]
For every \(j\in J_x(C)\), the construction returns \(\mathcal M_j=(Q_j,H_j)\in\mathfrak F_x(C,\varepsilon)\) with \(s_{\mathcal M_j}(x)=j\). Conditional on an exact input matrix \(C\), it computes \(J_x(C)\), all ratio frontiers, and all at most \(n\) dense witness pairs \((Q_j,H_j)\) in \(O(n^3)\) exact real-arithmetic operations. This order is optimal for dense all-witness output in the worst case. If \(C\) is rational, every returned \(Q_j\) and \(H_j\) is rational.
\end{lemma}
\begin{proof}
Put \(A=C_{-x,:}\). The \(p\) rows of \(C\) are linearly independent by the displayed rank hypothesis, so the \(p-1\) rows of \(A\) are linearly independent and
\[
 \operatorname{rank}(A)=p-1,
 \qquad
 \dim\ker(A)=n-p+1. \tag{1}
\]
The observed rows of \(T_0\) equal \(C\), and \(T_0W_0=I_n\). Hence, for every \(i\in\mathsf O\setminus\{x\}\) and every \(k\in\{x\}\cup\mathsf L\),
\[
 (A(W_0)_{:,k})_i=(T_0W_0)_{ik}=0. \tag{2}
\]
Thus the columns of \(K_x=(W_0)_{:,\{x\}\cup\mathsf L}\) lie in \(\ker(A)\). They are linearly independent because they are a subfamily of the columns of the invertible matrix \(W_0\), and their number is \(1+(n-p)=n-p+1\). Equation (1) therefore proves that they form a basis of \(\ker(A)\).

For a fixed \(j\), write \(a_j\) for column \(j\) of \(A\). If \(\operatorname{rank}(A_{:,-j})=p-1\), its columns span \(\mathbb R^{p-1}\), so there is a vector \(\beta\) with \(a_j=A_{:,-j}\beta\). The vector \(z\) defined by \(z_j=1\) and \(z_{-j}=-\beta\) then satisfies \(Az=0\) and \(z_j\ne0\). Conversely, if \(Az=0\) and \(z_j\ne0\), then
\[
 a_j=-\sum_{k\ne j}\frac{z_k}{z_j}a_k. \tag{3}
\]
Since \(A\) has rank \(p-1\), equation (3) shows that deleting \(a_j\) does not change its column span, and hence \(\operatorname{rank}(A_{:,-j})=p-1\). Because the columns of \(K_x\) span \(\ker(A)\), a vector \(z\in\ker(A)\) with \(z_j\ne0\) exists exactly when row \(j\) of \(K_x\) is nonzero. Combining this equivalence with the definition of \(J_x(C)\) proves the displayed admissibility test.

We next prove that all matchings can be obtained from one graph search. Since \(T_0=W_0^{-1}\), the cofactor formula gives, for every \(j\),
\[
 C_{xj}=(T_0)_{xj}
 =(-1)^{x+j}\frac{\det((W_0)_{-j,-x})}{\det(W_0)}. \tag{4}
\]
For \(j\in J_x(C)\), the left side and the denominator in (4) are nonzero. Thus \(\det((W_0)_{-j,-x})\ne0\). A nonzero term in the Leibniz expansion supplies a perfect matching of the nonzero pattern of this square submatrix. Equivalently, in the bipartite graph with all \(n\) row vertices and the \(n-1\) columns other than \(x\), there is a matching saturating all columns and leaving row \(j\) unmatched.

Let \(M_0\) be the single column-saturating matching computed by the algorithm and let \(j_0\) be its unique unmatched row. For any \(j\in J_x(C)\), let \(M_j^*\) be a column-saturating matching leaving \(j\) unmatched, whose existence was just proved. If \(j\ne j_0\), the symmetric difference \(M_0\mathbin\triangle M_j^*\) is a disjoint union of alternating cycles and one alternating path from \(j_0\) to \(j\). Consequently \(j\) is reached by the alternating search rooted at \(j_0\). Flipping \(M_0\) along the stored search-tree path to \(j\) matches \(j_0\), unmatches \(j\), and leaves every column matched exactly once. For \(j=j_0\), the empty path leaves \(M_0\) unchanged. This proves the validity of every shared-graph matching used by the construction.

It remains to verify the shear and structural identities. Because \(j\in J_x(C)\), row \(j\) of \(K_x\) is nonzero. If \((W_0)_{jx}=0\), there is therefore an \(\ell\in\mathsf L\) with \((W_0)_{j\ell}\ne0\). Since \(\ell\ne x\), the matrix \(E=e_\ell e_x^\top\) satisfies \(E^2=0\), and hence
\[
 T_j=(I_n+E)T_0,
 \qquad
 W_j=T_j^{-1}=W_0(I_n-E). \tag{5}
\]
The multiplication on the right in (5) changes only column \(x\):
\[
 (W_j)_{:,k}=(W_0)_{:,k}\quad(k\ne x),
 \qquad
 (W_j)_{jx}=(W_0)_{jx}-(W_0)_{j\ell}=-(W_0)_{j\ell}\ne0. \tag{6}
\]
The multiplication on the left defining \(T_j\) changes only latent row \(\ell\), so \((T_j)_{\mathsf O,:}=C\). If no shear is required, these same conclusions hold with \(T_j=T_0\) and \(W_j=W_0\). In particular, the shared matching uses unchanged entries of \((W_j)_{:,-x}\), and its added edge from row \(j\) to column \(x\) is nonzero by (6).

Let \(\Pi_j\) be oriented as in the construction. Every diagonal entry of \(\Pi_jW_j\) is the nonzero entry selected by the augmented matching, so \(D_j\) is invertible. Therefore
\[
 Q_j=D_j^{-1}\Pi_jW_j,
 \qquad
 H_j=D_j^{-1}\Pi_j   \tag{7}
\]
are invertible, \(H_j\) is monomial, and the diagonal of \(Q_j\) is one. Direct multiplication gives
\[
 Q_j^{-1}H_j
 =W_j^{-1}\Pi_j^\top D_jD_j^{-1}\Pi_j
 =T_j. \tag{8}
\]
Thus \([Q_j^{-1}H_j]_{\mathsf O,:}=C\). Because the augmented matching assigns row \(j\) to column \(x\), row \(x\) of \(\Pi_j\) is \(e_j^\top\), whence row \(x\) of \(H_j\) has its unique nonzero entry in column \(j\). This proves \(s_{\mathcal M_j}(x)=j\). Moreover,
\[
 (Q_j^{-1})_{xx}
 =(T_j\Pi_j^\top D_j)_{xx}
 =(T_j)_{xj}(D_j)_{xx}
 =C_{xj}(D_j)_{xx}\ne0. \tag{9}
\]
For an invertible matrix, cofactor expansion of the \((x,x)\) entry of its inverse yields
\[
 \det((Q_j)_{-x,-x})=\det(Q_j)(Q_j^{-1})_{xx}\ne0. \tag{10}
\]
Equations (7)--(10), together with the unit diagonal and monomial properties, verify every membership condition in \(\mathfrak F_x(C,\varepsilon)\).

We finally count arithmetic operations. Gaussian elimination can extend the \(p\) independent rows of \(C\) to \(T_0\) and invert \(T_0\) in \(O(n^3)\) operations. Reading all rows of \(K_x\) and all treatment loadings costs \(O(n^2)\). A standard augmenting-path implementation obtains the one column-saturating matching by at most \(n-1\) augmentations, each scanning at most \(O(n^2)\) edges, hence in \(O(n^3)\) operations; the alternating-reachability forest and all path flips cost \(O(n^2)\) in total. For a fixed \(j\), the shear changes one row of \(T_0\) and one column of \(W_0\), while forming the dense matrices in (7) costs \(O(n^2)\). There are at most \(n\) admissible indices, so dense witness materialization costs \(O(n^3)\). Computing all vector ratios costs \(O(np)\), and even pairwise exact deduplication of at most \(n\) vectors of length at most \(p\le n\) costs \(O(n^3)\). The entire computation is therefore \(O(n^3)\).

This rate matches worst-case dense output size. Choose distinct positive rational numbers \(t_1,\ldots,t_n\), put \(C_{xj}=1\), and let the other \(p-1\) rows be \(t_j,t_j^2,\ldots,t_j^{p-1}\). The usual Vandermonde determinant formula shows that \(C\) has row rank \(p\), while after deleting row \(x\) and any one column, any \(p-1\) remaining columns form a nonsingular matrix after factoring one nonzero \(t_j\) from each selected column. Thus \(J_x(C)=\mathsf S\), and the response vectors are distinct because their coordinate in the row carrying \(t_j\) is \(t_j\). Materializing one witness for each of these \(n\) values as a pair of \(n\)-by-\(n\) arrays has \(\Theta(n^3)\) scalar slots. Finally, if \(C\) is rational, Gaussian elimination may choose a rational basis extension; inversion, the unit shear, row permutation, and division by the nonzero diagonal entries in (7) preserve rationality. This completes all claims.
\end{proof}

\begin{lemma}[lem:canonical-exact-c-law-witness]\label{lem:canonical-exact-c-law-witness}
For every \(j\in J_x(C)\), let \(\mathcal M_j=(Q_j,H_j)\) be the fixed-law completion returned by \(\operatorname{Cert}(C,x)\), and set
\[
 \widetilde{\mathcal M}_j=(n,Q_j,H_j,\varepsilon).
\]
Then \(\widetilde{\mathcal M}_j\in\mathfrak F_x^{\mathrm{law}}(P_X)\), \(\bar s_{\widetilde{\mathcal M}_j}(x)=j\), and
\[
 \rho_x(\widetilde{\mathcal M}_j)
 =\left(\frac{C_{ij}}{C_{xj}}\right)_{i\in\mathsf O\setminus\{x\}}.
\]
If \(C\) is rational, \(Q_j\) and \(H_j\) are rational.
\end{lemma}
\begin{proof}
Fix \(j\in J_x(C)\).  By \(\textnormal{lem:shear-matching-sufficiency}\), the certificate returns \(\mathcal M_j=(Q_j,H_j)\in\mathfrak F_x(C,\varepsilon)\), satisfies \(s_{\mathcal M_j}(x)=j\), and has rational \(Q_j,H_j\) whenever \(C\) is rational.  Membership in the fixed-law fiber and \(\textnormal{def:completion-fiber}\) give
\[
 Q_{j,ii}=1,\qquad
 H_j\ \text{monomial},\qquad
 \det(Q_j)\ne0,\qquad
 \det((Q_j)_{-x,-x})\ne0,\qquad
 [Q_j^{-1}H_j]_{\mathsf O,:}=C.                                            \tag{12}
\]

Consider the canonical embedding
\[
 \widetilde{\mathcal M}_j=(n,Q_j,H_j,\varepsilon).
\]
Its observed mixing matrix, as defined in \(\textnormal{def:observational-law-completion}\), is \(C\) by (12), and therefore its observed law is
\[
 \mathcal L\!\left([Q_j^{-1}H_j]_{\mathsf O,:}\varepsilon\right)
 =\mathcal L(C\varepsilon)=P_X.                                            \tag{13}
\]
The independent, nondegenerate, and non-Gaussian source requirements in \(\textnormal{def:law-completion-fiber}\) follow respectively from \(\textnormal{ass:independent-sources}\), \(\textnormal{ass:nondegenerate-sources}\), and \(\textnormal{ass:nongaussian-sources}\).  Its no-zero-column and no-proportional-column requirements follow from \(\textnormal{ass:nonzero-columns}\) and \(\textnormal{ass:distinct-directions}\), because the observed mixing matrix is exactly \(C\).  The remaining diagonal, monomial, observational-solvability, and postintervention-solvability requirements are the first four facts in (12).  Thus
\[
 \widetilde{\mathcal M}_j\in\mathfrak F_x^{\mathrm{law}}(P_X).              \tag{14}
\]

The unique nonzero entry of row \(x\) of \(H_j\) is in column \(j\).  Since \(H_j\) is monomial, column \(j\) has no other nonzero entry.  Writing \(h=(H_j)_{xj}\ne0\), equation (12) consequently gives, for every \(i\in\mathsf O\),
\[
 C_{ij}=(Q_j^{-1}H_j)_{ij}=(Q_j^{-1})_{ix}h.                               \tag{15}
\]
At \(i=x\), \(\textnormal{def:admissible-source-set}\) gives \(C_{xj}\ne0\), so (15) also verifies \((Q_j^{-1})_{xx}\ne0\).  Division of (15) by its \(i=x\) instance is therefore legitimate and yields
\[
 \frac{(Q_j^{-1})_{ix}}{(Q_j^{-1})_{xx}}
 =\frac{C_{ij}}{C_{xj}}
 \qquad(i\in\mathsf O\setminus\{x\}).                                      \tag{16}
\]
This is the coordinatewise algebra underlying \(\textnormal{lem:intervention-ratio}\).  Applying the response formula \(\textnormal{lem:law-hard-intervention-response}\) to (14) and then using (16) proves
\[
 \rho_x(\widetilde{\mathcal M}_j)
 =\left(\frac{C_{ij}}{C_{xj}}\right)_{i\in\mathsf O\setminus\{x\}}.          \tag{17}
\]
Finally, \(\textnormal{def:law-assigned-source}\) and (12) give
\(s_{\widetilde{\mathcal M}_j}^{\mathrm{law}}(x)=j\).  The relevant observed mixing column is \(c_j\), and \(\textnormal{ass:distinct-directions}\) makes \(j\) the unique reference direction with that span.  Hence \(\textnormal{def:law-aligned-source}\) gives
\[
 \bar s_{\widetilde{\mathcal M}_j}(x)=j.
\]
The rationality assertion is inherited directly from \(\textnormal{lem:shear-matching-sufficiency}\).
\end{proof}

\section{Interpretation}

The causal object is the set of equilibrium responses over the same-law structural-completion fiber, not the rank-and-ratio formula by definition. Full-law OICA uniqueness first identifies the unordered projective directions; the necessity and constructive-sufficiency proof then shows that deletion rank selects exactly the directions attainable as the intervened equation's source. Representative invariance makes the resulting finite set an observational-law functional. Exact population cumulants evaluate that functional pointwise under all finite moments, whereas finite-sample root-\(N\) uncertainty requires the separate common-order quantitative margins.

\section*{Technical internal limitation (diagnostic only)}

Source minimality means exactly one distinct disturbance per structural equation, with no source sharing or splitting, and the exact-representative construction varies latent structural rows while holding one irreducible OICA direction set fixed. Enlarging the fiber to graphically irreducible completions, duplicated or split sources, or stability-restricted dynamics can change the frontier and is not covered. The heterogeneous-order evaluator assumes known source count, exact population cumulants, and finite moments of every order; it is pointwise and has no uniform stopping order. Moment determinacy is not used positively, although the nonuniformity witnesses can be chosen moment determinate. The common-order iid theorem retains only its fixed-order source-moment envelope and quantitative tensor, loading, denominator, and deletion-rank margins. The fourteen-stock Hong Kong series is a consumer of the population certificate only: its serial dependence and mechanism changes require a separate dependent-time-series theorem and an empirical separation certificate.

\section{Honest scope}

The paper delivers the scalar and vector law-fiber frontiers, constructive sharpness, exact same-law countermodels, representative invariance, and conditional worst-case output-optimal \(O(n^3)\) all-witness enumeration from an exact irreducible mixing representative. With known \(n\), exact population cumulants, normalized full-row-rank pairwise-nonproportional mixing, independent nondegenerate non-Gaussian sources, and all finite source moments, it also gives a terminating representation-dependent multi-order direction evaluator. It does not claim a uniform cumulant cutoff or a heterogeneous-order finite-sample rate; impossibility holds even for moment-determinate near-Gaussian witnesses. Simultaneous root-\(N\) inference remains restricted to the nonempty separated common-order iid class with its fixed-order moment envelope. Stability-refined completions, graphically irreducible fibers without source minimality, common-completion joint-intervention matrices, and dependent-time-series inference for the Hong Kong data are outside scope.

\begin{thebibliography}{99}

  \bibitem{DaiAlbrechtSpirtesZhang2026} Haoyue Dai, Immanuel Albrecht, Peter Spirtes, and Kun Zhang (2026), Distributional Equivalence in Linear Non-Gaussian Latent-Variable Cyclic Causal Models: Characterization and Learning, International Conference on Learning Representations, arXiv:2603.04780.

  \bibitem{TramontanoKivvaSalehkaleybarDrtonKiyavash2024} Daniele Tramontano, Yaroslav Kivva, Saber Salehkaleybar, Mathias Drton, and Negar Kiyavash (2024), Causal Effect Identification in LiNGAM Models with Latent Confounders, International Conference on Machine Learning, arXiv:2406.02049.

  \bibitem{TramontanoDrtonEtesami2026} Daniele Tramontano, Mathias Drton, and Jalal Etesami (2026), Parameter Identification in Linear Non-Gaussian Causal Models under General Confounding, arXiv:2405.20856, version 2.

  \bibitem{LacerdaSpirtesRamseyHoyer2008} Gustavo Lacerda, Peter Spirtes, Joseph Ramsey, and Patrik O. Hoyer (2008), Discovering Cyclic Causal Models by Independent Components Analysis, Proceedings of the Twenty-Fourth Conference on Uncertainty in Artificial Intelligence, arXiv:1206.3273.

  \bibitem{BongersForrePetersMooij2021} Stephan Bongers, Patrick Forré, Jonas Peters, and Joris M. Mooij (2021), Foundations of Structural Causal Models with Cycles and Latent Variables, Annals of Statistics 49, 2885--2915.

  \bibitem{WangSeigal2026} Kexin Wang and Anna Seigal (2026), Identifiability of Overcomplete Independent Component Analysis, Bernoulli, DOI 10.3150/25-BEJ1929, arXiv:2401.14709.

  \bibitem{Comon1994} Pierre Comon (1994), Independent Component Analysis, A New Concept?, Signal Processing 36, 287--314.

  \bibitem{ShohatTamarkin1943} James A. Shohat and Jacob D. Tamarkin (1943), The Problem of Moments, American Mathematical Society.

  \bibitem{ChenBickel2006} Aiyou Chen and Peter J. Bickel (2006), Efficient Independent Component Analysis, Annals of Statistics 34, 2825--2855.

  \bibitem{AnandkumarGeJanzamin2015} Anima Anandkumar, Rong Ge, and Majid Janzamin (2015), Learning Overcomplete Latent Variable Models through Tensor Methods, Conference on Learning Theory, arXiv:1408.0553.

  \bibitem{PodosinnikovaEtAl2019} Anastasia Podosinnikova, Amelia Perry, Alexander S. Wein, Francis Bach, Alexandre d'Aspremont, and David Sontag (2019), Overcomplete Independent Component Analysis via SDP, Proceedings of the Twenty-Second International Conference on Artificial Intelligence and Statistics, arXiv:1901.08334.

  \bibitem{VirtaNordhausenOja2016} Joni Virta, Klaus Nordhausen, and Hannu Oja (2016), Joint Use of Third and Fourth Cumulants in Independent Component Analysis, arXiv:1505.02613.

  \bibitem{NandyMaathuisRichardson2017} Preetam Nandy, Marloes H. Maathuis, and Thomas S. Richardson (2017), Estimating the Effect of Joint Interventions from Observational Data in Sparse High-Dimensional Settings, Annals of Statistics 45, 647--674.

  \bibitem{LeeShenTruong2021} Seonjoo Lee, Haipeng Shen, and Young K. Truong (2021), Sampling Properties of Color Independent Component Analysis, Journal of Multivariate Analysis 182, article 104692.

  \bibitem{HuangEtAl2020} Biwei Huang, Kun Zhang, Jiji Zhang, Joseph Ramsey, Ruben Sanchez-Romero, Clark Glymour, and Bernhard Schölkopf (2020), Causal Discovery from Heterogeneous or Nonstationary Data, Journal of Machine Learning Research 21, 1--53.

  \bibitem{SchkodaRobevaDrton2024} Daniela Schkoda, Elina Robeva, and Mathias Drton (2024), Causal Discovery of Linear Non-Gaussian Causal Models with Unobserved Confounding, arXiv:2408.04907.

  \bibitem{DrtonEtAl2025} Mathias Drton, Marina Garrote-López, Niko Nikov, Elina Robeva, and Y. Samuel Wang (2025), Causal Discovery for Linear Non-Gaussian Models with Disjoint Cycles, Proceedings of the Forty-First Conference on Uncertainty in Artificial Intelligence, Proceedings of Machine Learning Research 286, 1064--1083, arXiv:2507.10767.

  \bibitem{Kruskal1977} Joseph B. Kruskal (1977), Three-way arrays: rank and uniqueness of trilinear decompositions, with application to arithmetic complexity and statistics, Linear Algebra and its Applications 18, 95--138, DOI 10.1016/0024-3795(77)90069-6.

  \bibitem{Marcinkiewicz1939} Jozef Marcinkiewicz (1939), Sur une propriete de la loi de Gauss, Mathematische Zeitschrift 44, 612--618, DOI 10.1007/BF01210677.

  \bibitem{GoyalVempalaXiao2014} Navin Goyal, Santosh Vempala, and Ying Xiao (2014), Fourier PCA and Robust Tensor Decomposition, Proceedings of the 46th Annual ACM Symposium on Theory of Computing, 584--593, arXiv:1306.5825v5.

  \bibitem{TramontanoKivvaSalehkaleybarDrtonKiyavash2025} Daniele Tramontano, Yaroslav Kivva, Saber Salehkaleybar, Mathias Drton, and Negar Kiyavash (2025), Causal Effect Identification in lvLiNGAM from Higher-Order Cumulants, International Conference on Machine Learning, arXiv:2506.05202v2.

\end{thebibliography}

\end{document}